%% PREAMBLE
%% -----------------------------------------------------
\documentclass[12pt]{article}
%https://www.overleaf.com/4412272425fgkjfcmkgwgs
\usepackage[utf8]{inputenc}
\usepackage{graphicx}
\usepackage{xcolor}
\usepackage{wrapfig}
\usepackage[margin=0.5in,includefoot]{geometry}
\usepackage{booktabs}
\usepackage{lmodern,textcomp}
\usepackage{amsmath}
\usepackage{amsthm}
\usepackage{amssymb}
\usepackage{fullpage}
\usepackage[english]{babel}
\usepackage{breqn}
%\usepackage[section]{placeins}
%\usepackage[subsection]{placeins}
\usepackage{placeins}
\usepackage{natbib}
\usepackage{appendix}
\usepackage{footnote}
\makesavenoteenv{tabular}
\makesavenoteenv{table}
\usepackage{fancyvrb}
\usepackage{varwidth}
\usepackage{verbatimbox}
\usepackage{hyperref}
\usepackage[normalsize,justification=centering,singlelinecheck=off]{caption}
\usepackage[singlespacing]{setspace}
\theoremstyle{definition}
\newtheorem{definition}{Definition}
\theoremstyle{remark}
\newtheorem*{remark}{Remark}
\renewcommand{\baselinestretch}{1.25}
% \usepackage[dvipsnames]{xcolor}

\usepackage{abstract}

% Modify the abstract width
\renewcommand{\absleftindent}{0pt}
\renewcommand{\absrightindent}{0pt}

\definecolor{mygreen}{RGB}{60,130,80}
\definecolor{uv}{RGB}{8, 153, 158}

%% TEXT
%% -----------------------------------------------------

\title{\vspace{-2cm}
So alike, yet so different:\\
A non-parametric variance decomposition of \\ 
unstable employment cycles\thanks{We would like to thank participants at the Spanish Economic Association Meeting, 
the  ``New Frontiers for Policy Evaluation: Micro and Macro Perspectives'' Conference in NYUAD, the 2025 Econometric Society World Congress, The 1st Aarhus Workshop on Labor Economics, the Workshop on Employment Dynamics and Wage Structures in the 21st Century: Bridging Microdata and Policy, and seminar participants at UCL,  University of Mannheim and UAB for very useful comments. 
%The data basis of this study is the factually anonymous Sample of Integrated Labour Market Biographies of the Institute for Employment Research (IAB) (version 1975-2017). 
We thank Rob Shimer and Imran Rasul  for insightful discussion.
Data access to the SIAB was provided via on-site use at the Research Data Centre (FDZ) of the German
Federal Employment Agency (BA) at the Institute for Employment Research (IAB) and
subsequently remote data access.}}

\author{Maia Güell\footnote{CUNEF Universidad,  University of Edinburgh, CEPR, FEDEA, and IZA. Email:  maia.guell@gmail.com},
Cristina Lafuente\footnote{University of Bath. Email: 
clm96@bath.ac.uk},
Jenny Peters\footnote{University of Edinburgh. Email: jenny.peters@ed.ac.uk},
Manuel Sánchez\footnote{RavenPack. Email: manuzrp@gmail.com}, 
H\'{e}l\`{e}ne Turon\footnote{ University of Bristol and IZA. Email: Helene.Turon-Lacarrieu@bristol.ac.uk}
}



\date{April 2026}

\begin{document}
\vspace{-1.5cm}

\maketitle

\vspace{-1.5cm}
\begin{abstract}
\singlespacing 

%This paper examines employment instability in Germany and Spain—two countries with distinct labor market institutions and macroeconomic outcomes but similar rates of unstable employment. We define employment instability as any period of non-regular employment, including unemployment and atypical work, before transitioning into stable (permanent) employment. This combined-unstable-state approach offers a broader perspective on income volatility than focusing on unemployment alone. A key difference is the prevalence of mini-jobs in Germany—low-pay, low-hour contracts with continuous but weak labor market attachment—versus temporary contracts in Spain, which offer higher wages but strict duration limits and abrupt terminations, resulting in extreme labor market attachment.
%Using large-scale administrative data, we apply an extended variance decomposition methodology (ABS method) that makes no functional form assumptions to disentangle the roles of individual heterogeneity, duration dependence, and aggregate conditions in transitions from unstable to stable employment.
%Despite similar survival rates of unstable employment, we find stark differences in underlying mechanisms. In Spain, duration dependence is positive, meaning the probability of exiting instability decreases over time, whereas in Germany, it is negative, suggesting an increasing likelihood of transition to stability. These findings challenge the notion that stronger attachment to atypical employment leads to greater employment stability.
%While the decomposition avoids strong functional form assumptions, it requires two completed spells per individual. The use of comprehensive administrative data helps mitigate concerns about sample selection, though we discuss its potential effects and how larger T reduces selection bias.


\end{abstract}

%\par
%\vspace{1cm}


\noindent \textbf{Keywords}: administrative social security data, unstable employment spells, temporary contracts, mini-jobs, duration dependence. \\
\textbf{JEL Classification}: {J21}, {J64}, {J80}.

\textcolor{magenta}{Check ``Data sources´´ in notes of each table}
\textcolor{mygreen}{H\'el\`ene's comments in green}

\textcolor{magenta}{The method implies sample selection, which affects heterogenity (underestime); we do many things and detect who we are missings. Most of the missings are not so much our intererest. 
Overall results 
not much affected, specially duration dependence which is our interest..... The two-spell requirement implies selection on fast-cyclers and understates heterogeneity, but leaves duration dependence, ou rkey object of interest, essentially unaffected in both countries.}
\newpage

\section{Motivation or Introduction}


Intro should include something about Institutional Setting.

\begin{figure}[h!]
    \centering
    \caption{Incidence of Atypical Employment}
    \vspace{-0.5cm}
\begin{center}
\includegraphics[width=0.8\textwidth]{imgs/atypical_incidence_by_year.png}
\end{center}
\label{emp_stock_atypical}
\end{figure}

\textcolor{mygreen}{Explain policy relevance of 1. Unstable cycles, 2. Duration dependence vs heterogeneity to feed into targeting active labour market policies and income support policies.}

\textcolor{mygreen}{Rationale for looking at unstable spells. Relate to literature on unemployment spells and duration. Rising share of intermediate state of unstable or short employment. Matters for cross-country comparison as many countries, particularly with strong EPL on permanent contracts, have seen as sharp rise in atypical contracts which can take several different forms, e.g. mini-jobs in Germany, temporary jobs in Spain and France, zero-hour contracts in the UK, which all have in common that they provide workers with low and volatile incomes for sometimes extended periods of time.}

\textcolor{mygreen}{Have a section on context to explain different contracts, unemployment rates, active labour market policies, long-term unstable spells incidence, income dynamics in the two countries?}


\section{Active labour market policies}

\subsection{Spain}
\subsubsection{`Unstable' employment}
2002-2006: Labour Market Insertion Contract: direct employment programme implemented primarily within the Spanish public sector, aimed at individuals who had difficulties entering the labour market, and introduced with the purpose of reducing unemployment. In 2006, Spanish authorities decided to revoke it. Empirical research found its outcomes were poor relative to standard temporary contracts.

2010 Zapatero reform. It introduced a new employment-promotion permanent contract with maximum severance pay of 33 days' wages per year of service. Economic reasons for fair dismissal were broadened to include current and expected losses or a persistent reduction in revenues.  The reform also introduced German-style government-subsidised short-time working (the ERTE scheme), where companies can reduce the working week and the state unemployment system compensates affected workers with partial unemployment payments. It also penalised the use of successive temporary contracts: after three years of temporary contracts, a company would be obliged to take on the worker permanently.

2012 Rajoy reform included a large cut in statutory dismissal costs for workers on permanent employment contracts; the ability of employers to impose unilateral changes to employment conditions or suspend contracts when a company had seen a reduction in revenue over two trimesters.

COVID 2020-21 prompted an unprecedented use of the ERTE (Expediente de Regulación Temporal de Empleo) job retention scheme. Job retention schemes (ERTE furlough schemes) prevented unemployment from sky-rocketing during the pandemic by allowing millions of workers to retain their jobs.

2021 Sanchez reform The general limit for successive temporary contracts was tightened — employees who in a 24-month period (previously 30 months) have been employed for more than 18 months (previously 24 months) will acquire the status of permanent employees. 

The 2021 reform has had measurable effects on the labour market. The temporary employment rate fell from an average of 29.7\% in the period 2014–2019 to 12.7\% in 2024. This reduction has occurred across various sectors, age groups and regions, and has led to greater employment stability, although job turnover has increased and the number of contracts registered has decreased. The labour reform, combined with minimum wage increases and other policies, not only reduced temporary employment from 26\% to 12.2\% but also increased the active population by one million people.

\subsubsection{Unemployed job search}
The central activation tool is the compromiso de actividad — a formal agreement signed upon claiming benefits. This agreement is a document signed when applying for unemployment. It includes activities aimed at active job searching, reintegration into the labour market, and improving employability. The public employment service is obligated to provide a personalised plan tailored to each person's profession, training, and experience, with progress monitoring by a tutor and the offering of activities and services that facilitate returning to the labour market. The unemployed person commits to participating in educational activities, accepting suitable job offers, and proving compliance with the agreed actions.

The penalty structure for failing activation requirements is graduated. Sanctions are classified as minor and major. Minor infractions — such as failing to meet activity requirements without justification — are sanctioned as follows: 1st infringement: loss of benefit for one month; 2nd: three months; 3rd: six months; 4th: permanent extinction of the benefit. Major infractions include refusing a suitable job offer and refusing to participate in programmes proposed by the SEPE to improve employability: 1st infringement: loss of benefit for three months; 2nd: six months.

Several targeted programmes have been created for those who exhaust standard benefits:
Renta Activa de Inserción (RAI): An active insertion benefit for those with no other entitlements, targeting particularly vulnerable groups. 
PREPARA Programme (2011–2016): A post-crisis programme combining a small income supplement with personalised guidance. 
Activation Programme for Employment — PAE (2014 onwards): The PAE, agreed between the government and the social partners in December 2014, aimed to improve the employability of the long-term unemployed in situations of special need. 

\subsection{Germany}
\subsubsection{`Unstable' employment}
The 2001 Part-Time and Fixed-Term Employment Act (TzBfG)
Just before the Hartz era, the government overhauled the legal framework for non-standard contracts. In November 2000, the German parliament passed a new Act on part-time work and fixed-term employment relationships (Gesetz über Teilzeitarbeit und befristete Arbeitsverträge). fixed-term contracts require an objective reason (covering for another worker, for example), though fixed terms of up to two years without an objective reason are permitted and may be extended up to three times.  

Hartz I \& II (2003): Liberalising Temporary and Marginal Work
The first three stages of the reforms (Hartz I–III) sought to improve job search efficiency and employment flexibility.
Temporary agency work was substantially deregulated. Prior to Hartz, maximum assignment durations had been progressively lengthened — from three months originally to twelve months by 1997. The Hartz reforms, implemented between 2004–2005, then allowed for unlimited assignment durations under specific conditions, thereby greatly increasing the flexibility of temporary employment. 

Mini-jobs and midi-jobs were expanded. Mini-jobs existed before the Hartz laws, but Hartz II raised the monthly salary cap from €325 to €400 and created a higher tier called midi-jobs — both were low-wage contracts for a small number of hours, with lower or gradually rising taxes and social insurance payments. 

Kurzarbeit: Internal Flexibility as an Alternative to Lay-offs: a social insurance programme whereby employers reduce their employees' working hours instead of laying them off; the government normally provides an income replacement rate of 60\% (more for workers with children) for each hour not worked. During COVID-19 the scheme was expanded to cover over 10 million workers at its peak. 

\subsubsection{Unemployed job search}

The Two-Tier Benefit System and ALG I vs. ALG II Post-Hartz, Germany has two distinct benefit tracks with different activation conditions. Arbeitslosengeld I (ALG I) is contribution-based, for those who have paid into the unemployment insurance fund; Arbeitslosengeld II (ALG II, or "Hartz IV") is a needs-based benefit for those who have not contributed enough or whose ALG I is insufficient to cover living expenses. Hartz IV merged the previously separate long-term unemployment assistance and social welfare into a single flat-rate benefit, managed by the new Jobcentres — a critical institutional innovation.

The Hartz IV reform merged the federal-level unemployment agency with local-level welfare administration. The plan's objective was to reduce caseloads from 400 unemployed persons per agent to no more than 75 (for those aged 25 or under) or 150 (for those over 25). 

The Integration Agreement (Eingliederungsvereinbarung) The central tool for job search activation was the integration agreement. These documents contain common goals and mandatory efforts of the unemployed as well as services of the institution — like job offers, participation in labour market schemes, or further training. The unemployed are no longer regarded as passive recipients of benefits but are supposed to shape their individual process of integration together with their partners in labour offices and job centres.

Sanctions Under Hartz IV An unemployed Hartz IV recipient could see their benefit reduced by 10\% for three months if they did not attend a Jobcentre meeting, and by 30\% for three months if they refused a suitable job offer, occupational training, or interrupted a training programme already underway. Second offenders within a year lost 60\% of their benefits; third-time offenders lost all benefits. 

The Bürgergeld Reform (2023) The integration agreement was replaced by a softer cooperation plan (Kooperationsplan). After examining with a case officer what types of support are needed, a cooperation plan is created, which jointly defines concrete steps towards a new job. Both the claimant and the Jobcentre commit to the plan, and if there is disagreement, a mediation process can be initiated. 

Training Programmes and the "One-Euro Job"Beyond job placement, Germany has operated a wide range of ALMP training instruments. One controversial post-Hartz creation was the "one-euro job" (Arbeitsgelegenheit), a work-opportunity scheme paying €1 per hour on top of benefits: over 300,000 people were in one-euro jobs at a time, which paid an hourly rate of €1 in addition to ALG II benefits. Evaluations of these were mixed; they were primarily aimed at preserving work habits among the very long-term unemployed rather than labour market integration. 




\section{Comments UAB}

\textcolor{magenta}{HERE COMMENTS}

\begin{enumerate}
    \item Say upfront, the shareunempl our sample represents... Not sure wrt what (LF?), but people keep asking how weird is our sample. More explicit that PC do lose their jobs. we could do this, but we should decide if we want it with respect to everyone or only those who every experienced unstable
     \item Say minijobs: continuos-low-attach. It is in the slides, include here. Downplay role of having UI at the same time.
    \item  ⁠Show demographics including bin 0 and 1 wrt our sample. we have this information ready if needed
\item LITERATURE: see new papers. we should look at the papers claude found for us. Cite Erosa et al (2025), MCLV-recession

\item COUNTRIES and CONTRACTS:⁠  ⁠resample without the unempl. do we need this if we already do the temp only/mini only? also, we will lose a lot of people if we exclude the unemployed, especially in Germany


    
\end{enumerate}

\textcolor{mygreen}{Would it flow better to have methods first, then data, then descriptives?}


\section{Data and definitions}\label{sec:data}

For each country we use their respective administrative social security records.  These data track daily one person-spells 
for their working histories. These go back to when individuals first entered the labour market up to their most recent spell.
The very nature of these data makes them comparable, although we need to take into account the population covered might differ in each country.\footnote{
To make the samples in the two countries as comparable as possible, we exclude those with at least half of any employment spell in self-employment and civil servants.
In Spain we also exclude 
those with a permanent-discountinuous contract. For more details, see \cite*{sodiferentsoalike}.}
%Similarly, we take into account that differences in  labour market institutions can affect some classifications in the data in each country. 
In what follows, we explain in more detail the data sets for each country. 



The administrative dataset from the social security records of Germany is the \emph{Sample of Integrated Labour Market Biographies} (\cite*{SIAB}), SIAB hereafter, provided by the Institute for Employment Research (IAB), which started in 1975.\footnote{Observations for East Germany are available after 1991.} 
The SIAB is a yearly 2\% random sample of the workforce in Germany drawn from the Integrated Employment Biographies (IEB), 
which contains administrative records for the entire German workforce.\footnote{Given the rules of the social security notification procedure, daily employment records are reported annually, which means that the starting and ending year are always identical (there is an obligation of the employer to submit annual employment notifications). A continuous employment spell may therefore be distributed across several notifications.} 
Crucial for this paper, 
mini-jobs can be observed since April 1999 and  temporary and permanent 
contracts can be distinguished in the data since November 2011. Hence in this paper, we use the years 2011-2021.\footnote{These years include different phases of the business cycle as well as Covid. One of our robustness checks controls for this.} 

The administrative dataset from the social security records of Spain is the \emph{Muestra Continua de Vidas Laborales} (\cite*{MCVL}), MCVL hereafter.\footnote{Spanish Continuous Sample of Working Lives.} 
The MCVL is a yearly 4\% random sample of the workforce in Spain.
Crucial for this paper, temporary and permanent contracts can be
distinguished in the data since 1996. We will be able to analyse the same sample period as Germany, 2011-2021, as well as a longer period for robustness checks.

\textcolor{magenta}{Check data citation for Spain} 
example published with MCVL:
https://doi.org/10.5281/zenodo.6463315.

citation MCVL: see below
%Ministerio de Inclusión, Seguridad Social y Migraciones (2018) ‘Muestra Continua de Vidas Laborales’, accessed on 21 May 2018.



% \paragraph{Labour market states}
% The labour market states in the German data are categorized as follows: employed either full-time or part-time and holding a permanent or temporary contract. Among those in part-time job, the data sets distinguishes between those who are in marginal employment (i.e., a mini-job) and those who are not. The other states are: unemployed (receiving benefits or, if not, registered as a job-seeker in an employment agency),\footnote{Registered job-seekers information is available since 1997 with missing information for years 2005 \& 2006.} and in non-employment. The latter includes: spells  without any type of notifications, and registered in an employment agency as `not unemployed'. {Workers in partial retirement can be identified (but not those fully retired).}
% Any other state (such as self-employment) will appear as `blank' spell. After the comparability criteria described below, all blank spells present in the data prior to a stable employment will be considered as non-employment.\\

% The labour market states in the Spanish data are categorized as follows: employed with a permanent contract, employed with a temporary contract,
% unemployed (receiving benefits),
% self-employed and retired (receiving a pension). Any other state (such as non-employment or unemployed without receiving unemployment insurance) will appear as `blank' spell.\footnote{See \cite{lafuente2019unemployment} for more details.} 
% After the comparability criteria described below, all blank spells present in the data prior to a stable employment will be considered as non-employment.

% \paragraph{Measurement unit}
% Our unit of analysis are individual spells. Our samples for the two countries will include
% the following {completed and uncompleted spells:
% employment spells (in any type of contract)}, unemployment spells  as well as any spell of non-employment in between employment spells.
%  The idea is to capture workers that might be provisionally out of the labour force but are still attached to the labour market. At the same time, we do not want to simply capture instability due to early career stages and, therefore, we focus on workers who are at least 25 years old.\footnote{In a recent paper, \cite*{Jimenoetal2021} discuss employment for workers below 25 years in Spain during the last two recessions.} As explained below, we also restrict workers' age to be 60 years old.\\

% In order to make the two data sets as comparable as possible we  take the following steps:
% \begin{itemize}
% \item[(i)]{
% In the Spanish data, by design, workers included in the data start with an employment spell. While in the German data 
% it is possible to observe entry into the labour force, whether in employment or not. For comparability, in the German data set, we consider workers from the start date of their first employment spell.\footnote{In subsection \ref{sec:unstable_spells_def}, we discuss the possible under-reporting of duration for \emph{unstable spell}.}
% }

% \item[(ii)] {The German data set does not include spells of any civil servants.\footnote{See \cite*{Umkehrer2019}.} We therefore exclude this category in the Spanish data set.\footnote{Most civil servants do not appear in the Spanish data set as they are covered by a separate social security system. However, some employees of public firms and casual public workers sometimes do appear in the data. These are the observations we drop from the sample.} 
% Moreover, we also exclude employment spells in ``permanent discontinuous" contracts since these do not exist in Germany.\footnote{As mentioned before, these kind of contracts are targeted at workers of highly seasonal industries, like vegetable canning.}}
    
% \item[(iii)] We exclude in the German data workers in apprenticeships and student trainees, {given that in Spain until very recently these type of contracts and active labour market policies were not as significant as in Germany.}

% \item[(iv)] {Retired workers can be identified in the Spanish data but not in the German data (only if it is partial retirement). In Figure \ref{fig:Inactive} we explore the  inactivity rates by age group in Germany and Spain in 2021 (for previous years this Figure looks very similar). For both countries, at age 60 at least 50\% of the population is inactive. In order to have comparable samples in which retired workers are excluded, the maximum age allowed in our sample is 60. 
% }

% \begin{figure}[!h]
%     \caption{Inactivity Rates}
%     \label{fig:Inactive}
% \begin{minipage}[t]{1\textwidth}{}
%           \centerline{
%           \includegraphics[width=1.1\textwidth]{imgs/inactivity_rates.png}}
%       \end{minipage}
%   \\
% \centerline{\footnotesize{%Notes:  This figure  displays the inactivity rates by country.
% Source: \citet{Eurostat_inactivity} }}
% %Inactive population as a percentage of the total population, by sex and age (%) [LFSQ_IPGA__custom_1109587] https://ec.europa.eu/eurostat/databrowser/bookmark/e17f3a4b-7406-46b3-b480-975cece81c61?lang=en} 
% \end{figure}


% \item[(v)] {In the German data set, self-employment spells are not recorded.\footnote{See \cite*{Umkehrer2019}.} This means that a worker that only has self-employed spells in his/her working life will not be included in the data set. However, workers with working histories
% featuring some employment spell as a wage-earner as well as other spells of self-employment do appear in the data set (but these spells of self-employment will appear as `blank').
% To enhance comparability, in the Spanish data set we
% exclude workers with only self-employment spells in their working life. As for workers with mixed histories, in the Spanish data 
% we consider self-employment spells in between stable employment spells that are short enough\footnote{Less than 2 years, which is 40\% percent of the cases. This is also the maximum duration of a temporary contract. The distribution of durations of self-employment spells is very polarised, reflecting that this category covers very stable, almost managerial jobs and very short, contractor-like jobs. The later we consider unstable employment spells.} as non-employment to make it comparable to Germany. That is, these type of spells are considered unstable employment spells in the two countries. As for longer spells in these categories, we exclude them as they would not appear in the German data set.}

% \item[(vi)] Given the sample years used, in the German data set, uncompleted spells will have a maximum duration of 5 years. To enhance comparability, we make the same restriction in the Spanish data.

% \end{itemize}

% Our samples for the two countries will therefore include spells for: (i) employed workers (ii) unemployed workers (ii) non-employed in between employment spells (iv) self-employed (for short spells and Spain only). The inclusion of non-employment in (ii) is meant to capture workers who might be provisionally out of the labour force but are still attached to the labour market. Note that we include both complete and incomplete spells. Since our administrative data allows us to follow individuals since the beginning of their labour market activity, our spells are never left censored. They may however be right censored if the only unstable spell that we observe for a given individual is not complete by the end of the sample.\\

% Since our unit of observation is one spell per individual, the picture that we present departs from aggregate statistics on stocks: our analysis describes the individual experience of workers flowing into the ``unstable'' state (see definition below). Since we are giving the same weight to individuals with long spells as to individuals with short spells, we over-represent the latter relative to the former compared to an analysis of aggregate stocks. This explains, for example, why the fraction of spells in non-employment longer than one year is much smaller than in aggregate statistics.\footnote{That is, consider a quarterly labour force survey. A worker what has an unstable spell lasting for 2 quarters of more will appear in the stock-based data as unstable in both quarters. However, if a worker with a very short unstable spell in between quarters will not be captured by the labour force survey. Therefore, in our flow-based data, short spells are over-represented. This is because the frequency of administrative data is daily, and no such high-frequency survey data exists. This is explored in detail for Spain in \cite{lafuente2019unemployment}.}\\


% Other than the labour market status, both data sets include individual characteristics such as gender, age and education that we will use to undertake some heterogeneity analysis.

\subsection*{Spells: stable and unstable}\label{sec:unstable_spells_def}

Our approach is to define two states of the labour market: stable and unstable. A \emph{stable spell} is simply an employment spell 
with a permanent job. 
We define an \emph{unstable spell} as any combination of  atypical employment (temporary, mini-job and short self-employment) 
as well as non-employment. In other words, any labour market state that is not a permanent contract.\footnote{Part-time work is considered a stable state as long as it is a permanent contract, as it allows workers to be stable according to the hours worked.
}



We find this broad definition of \emph{unstable} useful for two reasons. First, atypical contracts are frequently observed with adjacent periods of non-employment. This combined-state allows us to include all forms of unstable labour market states.
Second, adopting a broad definition of instability allows for a cross-country comparison on the overall pattern of instability. 

The data is at the spell level.
A spell is defined as the uninterrupted time a person spends registered in a given labour state. Spells are measured in days. In the case of overlapping spells, or spells that end before the previous one starts, we apply the standard order of prevalence approach as described in \citet{perez2008muestra}. That is, a worker who simultaneously works in a permanent and a temporary job is classified as permanent for the duration of the permanent spell, as this is considered to be the main job. Similarly, employment spells of any kind take precedence over unemployment spells.\footnote{Workers receiving unemployment benefits during jobs usually take these in subsidised contracts, so we consider them employed.} In Spain, we use the expansion approach to unemployment as described in \citet{lafuente2019unemployment}. \textcolor{mygreen}{give some detail on this?} In Germany, the SIAB data allows us to differentiate directly between periods of inactivity and unemployment using the Jobseeker History data source.\footnote{Arbeitsuchendenhistorik - ASU / XASU, \cite*{fdzreport}.} \\


The empirical methodolody relies on completed spells. Therefore our definition of a \emph{completed  unstable spell} is as follows:

A completed unstable spell is any unstable spell that 
(i) starts during our sample time frame and (ii) ends in a stable spell. 
This unstable spell is the first one observed in our time frame. For the purpose of this paper, the term \emph{unstable spell} will be used to denote a \emph{completed} unstable spell.

The empirical methodology requires that we include two completed unstable spells per individual. Figure \ref{fig:tempdef} shows the sequence of spells for  individuals in our sample.
Since we only have two labour market states, unstable spells are preceded by a stable spell. 
The exception to this is for new entrants to the labour force, 
where we start counting their first unstable spell from their first atypical employment contract.
\textcolor{mygreen}{this new entrant rule is not very clear.}


\begin{figure}[h!]
    \centering
    \caption{Completed Spells Timeline in our sample}
    \vspace{-0.2cm}
    %, pure states}
    \begin{minipage}{1\textwidth} % choose width suitably
    \centerline{\includegraphics[width=0.8\linewidth]{imgs/spell_def_new.png}}
    \centering 
    \end{minipage}
    \label{fig:tempdef}
\end{figure}


%{Left-censored unstable spells that fall into our sample period are not considered.} 

Recall that unstable spells include unemployment. In fact, 
in both countries, the majority of transitions from stable to unstable are flows to unemployment and not to atypical contracts.\footnote{For Germany, among workers on a permanent contract (PC) only 20\% move to an atypical contract.}
\textcolor{magenta}{we need this number for spain.}

\textcolor{magenta}{Workers in our sample are not rare... }
\textcolor{mygreen}{In Data section, describe data breakdown by bins. Show time window – sample selection relationship. Show descriptive statistics for each (?) bin, or at least for data and sample. Characterise differences in terms of observables. Discuss potential differences in terms of unobservables. Announce sample selection exercise later in the paper.}

\textcolor{mygreen}{Document difference in income levels and income instability between unstable and stable state for each country?}

% \begin{figure}[h!]
%     \centering
%     \caption{Unstable Spell definition}
%     \vspace*{-0.2cm}
%     \begin{minipage}{1\textwidth} % choose width suitably
%     \centerline{\includegraphics[width=1\linewidth]{imgs/spell_definition.png}}
%     \footnotesize{Notes: Representation of a timeline of a worker. The time period encompassed by the bracket constitutes an `unstable spell'.}
%     \end{minipage}
%     \label{fig:tempdef}
% \end{figure}

%\begin{figure}[h!]
 %   \centering
  %  \caption{Completed Unstable Spells }
    %, pure states}
    %\begin{minipage}{1\textwidth} % choose width suitably
    %\centerline{\includegraphics[width=0.8\linewidth]{imgs/spell_definition_pure.png}}
  %  \footnotesize{Notes: Representation of a timeline of a worker with a completed unstable spell. The time period encompassed by the bracket constitutes an `unstable spell' in all three cases. }
   % \end{minipage}
  %  \label{fig:tempdef_pure}
%\end{figure}





\section{Methodology}\label{sec:method}

\textcolor{mygreen}{cite literature on duration dependence vs unobserved heterogeneity literature. Explain difficulty of disentangling the two but importance in terms of policy choices. Explain specificity and advantage of ABS method. Make clear the difference between individual duration distribution and population one.}

We use the variance decomposition method proposed by
\citet*{ABS}.
The starting point is to assume that
there are $N$ individuals with $J\geq2$ completed unstable spells
{in settings in which $N$ is large}. Our outcome of interest, unstable spell duration ($y$), is a random variable drawn independently for each individual $i$
from a probability distribution function $F_{i}(y)$, with mean $\mu_i$ and variance $\sigma^{2}_i$.

Let $y_{ij}$ be the duration for individual $i$ in spell $j$.
The population mean and variance of unstable spells are given, respectively, by


\begin{equation}
\mu_y = \frac{1}{N} \sum_{i=1}^{N} \int{ y_{ij} dF_{i}(y)} \hspace{0.4cm} \text{ and }  \hspace{0.4cm}
\sigma^{2}_y = \frac{1}{N} \sum_{i=1}^{N} \int{ (y_{ij}-\mu_{y})^{2} dF_{i}(y)},
\end{equation}


The population variance can be decomposed into two components: the \emph{within} and the \emph{between} components:

\begin{equation}\label{eq:main_decomp}
\sigma^2_{y} & = {\sigma^2_{b}}+ {\sigma^2_{w}}\\
 & ={\dfrac{1}{N}\sum_{i=1}^N (\mu_{i}-\mu_y)^2} + { \dfrac{1}{N}\sum_{i=1}^N \sigma^2_{i}} 
\end{equation}

That is, the between component is the variance in mean durations across individuals, and the within component is the mean of individual variances across the different spells of individuals.


The between component represents the variance that comes from heterogeneity in the distribution functions  $F_{i}(y)$ across individuals. Consider the following two extreme cases, as in \citet*{ABS}. First, all individuals draw from different, individual-specific distributions, but there are no differences within individuals (i.e. $\sigma^2_{i}=0$ for all $i$). In this case, the total variance in unemployment duration would be attributable to individual heterogeneity captured by
$\sigma^2_{b}$ in equation (\ref{eq:main_decomp}). 
Now consider the opposite case, in which all individuals draw from the same distribution function 
so all individuals have the same average duration (i.e.,
$\mu_{i}=\mu_{y}$ for all $i$). In this case, the variance in duration comes from
within an individual's draws in their different spells.

We now turn to the within component. 
Assume that the function $F(y)$ is given by a constant hazard rate $h$
of finding stable employment. Then the population and individual means
are $\mu_{y}=\mu_{i}=1/h$ and the standard deviations are $\sigma_{y}=\sigma_{i}=1/h$. In this case, by definition, there is no duration dependence 
and the total variance is given by a constant
for all individuals and all spells.
\citet*{ABS} label this component 
as the ``constant hazard within variance.'' 
We interpret this as a component coming from the
aggregate frictions in the labour market that prevent workers from finding a stable job immediately. The time that takes workers to find a stable job, or unstable spell duration, is drawn from an exponential distribution with parameter $h$. \textcolor{mygreen}{I don't think that's quite right. My understanding is that this component reflects the stochastic nature of the job finding process, i.e. even if everyone has the same job finding rate and it is constant with duration, individual durations will not all be the same because of different draws from a random variable distribution. Also, individuals experience very few draws from this distribution over their career, so this `risk' can't be averaged out at the individual level.}

However, to the extent that the hazard rate varies with elapsed duration in unstable spells (e.g., the longer in unemployment/temporary contracts, the longer it takes to leave these unstable cycles), then the variance will not be fully captured by the ``constant hazard within variance.'' 
Negative duration dependence manifests itself by increasing the variance relative to the mean of individual spells (\citealt*{StoyanDaley}). 

Duration dependence can be thought of as the distance between the variance and the squared mean, the within component from equation (\ref{eq:main_decomp})  can be expressed as:

\begin{equation}\label{eq:within_decomp}
 \sigma^2_{w}= & \dfrac{1}{N}\sum_{i=1}^N \sigma^2_{i}
  = & \dfrac{1}{N} \sum_{i=1}^N (\sigma^2_{i}-\mu^2_{i})+\dfrac{1}{N}\sum_{i=1}^N \mu^2_{i}=
  {\sigma^2_{e}+\sigma^2_{c}}
\end{equation}

The within component is formed by:  the excess variance, which an be positive, negative, or zero, and the constant hazard. 

In sum, we have that the total variance is given by the following three components: the between variance, the excess within variance and the constant within variance.

\begin{equation}\label{eq:main}
\sigma^2_{y} & ={\sigma^2_{b}} + { \sigma^2_{w}} & = {\sigma^2_{b}}+{\sigma^2_{e}+\sigma^2_{c}}\\   
\end{equation}


The three components correspond to heterogeneity, duration depedence and the constant. To visualize the contribution of these three components, we provide simulations that plot each of them in isolation, see Appendix \ref{app_simulations}. 


\paragraph{Estimation of the variance components}\label{sec: estimationABS}
Taking the
$N$ individuals with  $J\geq 2$ spells each,
ABS show that 
using only two completed spells per individual one can obtain unbiased estimators of the variance components with a large number of individuals.
In this case, 
the average and the variance of the two spells
are unbiased estimators of the individual means and variances.
And the sample mean and sample variance  
are unbiased estimators of the population mean and variance, where the sample
is composed by two spells from $n$ individuals.
These unbiased estimators of  $\mu_i$, $\sigma^{2}_i$,
$\mu_y$ and  $\sigma^{2}_y$ can be used to obtain unbiased estimators of the three components of the variance.

We have assumed so far that the two observations per individual are 
from the same individual distribution $F_{i}$. This may not necessarily
be the case for all applications, and $F_{i,1}$ may differ from $F_{i,2}$.
ABS show that the proposed estimate of the individual variance is unbiased
if and only if the two distributions have the same mean.
ABS argue that it is reasonable to assume this is the case if the mean of all first spells across individuals is the same as the mean of all second spells across individuals. In practice, to achieve this one has to first
regress the duration variable controlling for being the first or second spell and then do the decomposition on the residuals from the regression. This is represents
the variance that is not explained by the spell number. The same procedure can be applied for any other observable covariate.

To sum up, we have:

\begin{align}\label{eq:main2}
\hat{\sigma}^2_{y}  & = \hat{\sigma}^2_{s} + {\hat{\sigma}^2_{b}} + { \hat{\sigma}^2_{w}} \\
                    & = \hat{\sigma}^2_{s} + {\hat{\sigma}^2_{b}}+{\hat{\sigma}^2_{e}+\hat{\sigma}^2_{c}}\\   
\end{align}

where $\hat{\sigma}^2_{s}$ is the estimated variance explained by the spell number. Similary, if we added controls (e.g., controling for the business cycle), that would be captured by $\hat{\sigma}^2_{s}$ too.


For every decomposition, $\hat{\sigma}^2_{c} = 1.645$. This is because if the duration follows an exponential, then the duration in log days follows a log Weibull, or Gumbel, distribution. The variance of which equals $\pi^2 / 6 \approx 1.645$.
For more details, see Appendix \ref{baseline_var}.


% Given the two (completed) spell requirement, the ABS methodology 
% has limitations in dealing with  right-censored spells. We discuss the potential sample selection issues in our sample arising from this requirement below in section \ref{data} and Appendix \ref{appendix2spellreq}.


\textcolor{mygreen}{is there any role for measurement error? should we say something about that?}




\section{Descriptive statistics for Germany \& Spain} \label{sec: desc stats}

Given the data requirements by the method in section \ref{sec:method},  we sample from each country all those workers that fulfil the two spell requirement in our time frame. 
In practice,
our sample comprises individuals with two or more completed spells. As in \citet*{ABS}, we select the first two spells of each individual,  which minimizes issues with time-varying unobserved heterogeneity.

Table \ref{tab:desstats_dur} displays descriptive statistics for the duration of unstable spells for Germany and Spain. \textcolor{mygreen}{Table 1: is this for all workers who have experienced an unstable spell? What fraction of the population is this?}
Both distributions are right-skewed, displaying a relevant share of long durations. 
There are also some differences. 
On average, a worker in Germany spends 
374 days in an unstable spell, around a year; while in Spain is 912 days, almost 3 years. 
However, Germany is actually 
more dispersed than Spain, which is somewhat surprising given Spain's longer durations. 
The P90/P10 ratio is much higher in Germany than in Spain. Also, the coefficient of variation (std/mean) is higher in Germany at around 1.35 and in Spain around 1.05. 
Relative to its mean, the German distribution is more dispersed, 
with many short spells and a small share of  long ones. Despite the lower average  duration of unstable spells in Germany, the distribution shows more inequality. 

\textcolor{mygreen}{can we show a histogram of durations for each country and each spell type here? with a bar for spells that are censored because they are longer than the maximum window time. (and information in the table notes about the fraction of all censored spells.}

The comparison of the variance and 
log variance captures exactly this: in Germany, relative to Spain, the variance is smaller, but the log variance is higher. 
This is because the variance captures the absolute dispersion while the log variance captures the relative dispersion. \textcolor{mygreen}{This paragraph with variance and log variance is not clear. Footnote 10 is not needed.}\footnote
{Var$(g(Y)) \approx (g'(\mu))^2 \cdot \sigma^2$.
For $g(y) = log(y)$, we have $g'(y) = 1/y$, so $g'(\mu ) = 1/\mu $, giving  Var$(log(Y)) \approx \sigma^2 /\mu^2$.
This result follows from the Delta Method, \cite*{LehmannCasella1998}.} These descriptive statistics will map into our duration decomposition results. We will proceed with the variance decomposition in logs in the next section.



\begin{table}[!httb]
    \centering
    \caption{Descriptive Statistics: Duration unstable spells}
    \begin{tabular}{lll}
\toprule
 & Germany & Spain \\
\midrule
Mean (days) & 360.18 & 912.45 \\
Variance & 414,134.9 & 923,454.5 \\
Log (variance) & 2.568 & 1.855 \\
Min & 0 & 1 \\
p10 & 29 & 59 \\
p25 & 88 & 186 \\
p50 & 317 & 483 \\
p75 & 819 & 1,075 \\
p90 & 1,480 & 1,946 \\
Max & 6,884 & 10,265 \\
N & 81,768 & 69,406 \\[5pt]
\bottomrule
\end{tabular}

\label{tab:desstats_dur}
\end{table}   

Table \ref{tab:desstats_2+spell} provides descriptive statistics of our samples in the two countries. Both data sets have information on age, gender, nationality, and education. The samples are fairly comparable. We separate upper and lower secondary as there are many technical schools in Germany. We keep this distinction, but we find that these groups have a similar decomposition. 

\textcolor{mygreen}{could we have descriptive statistics of the underlying population here in Table 2? to answer the question: how weird is your sample.}

The table also displays the number of stable ($N_{S}$) and unstable spells  ($N_{U}$) for individuals in our sample, which can be more than two. 
Interestingly, the average $N_{U}$ is around two and the average $N_{S}$ is around three, which means that the two-completed spells that we study (see Figure \ref{fig:tempdef}) represent almost the whole working history of individuals in our sample.

\begin{table}[h!]
    \centering
    \caption{Descriptive Statistics}
\begin{tabular}{lcc}
\toprule
 & Germany & Spain \\
\midrule
Age & 36.60 & 41.35 \\
    & (7.17) & (7.37) \\
Female & 0.50 & 0.43 \\
    & (0.50) & (0.49) \\
Non-national & 0.26 & 0.18 \\
    & (0.44) & (0.38) \\
Primary & 5.81 & 12.16 \\
    & (0.00) & (0.00) \\
Lower Secondary & 20.10 & 39.62 \\
    & (0.00) & (0.00) \\
Upper Secondary & 53.55 & 26.22 \\
    & (0.00) & (0.00) \\
College & 20.54 & 22.00 \\
    & (0.00) & (0.00) \\
\midrule
$N_S$: number stable spells & 3.35 & 3.52 \\
   & (1.44) & (1.86) \\
$N_U$: number unstable spells & 2.85 & 2.50 \\
   & (1.35) & (1.37) \\
\hline \hline
$N$ & 81,768 & 69,406 \\
\bottomrule
\end{tabular}
 \label{tab:desstats_2+spell}
\end{table}

\textcolor{magenta}{Tables \ref{ind_DE} and \ref{ind_ES} in Appendix \ref{AppendixIndustry} show the distribution by industry of last employment for our sample.}
Despite some differences, employment in both countries is concentrated in services, especially trade and retail, transport, business-related activities, and social services. Manufacturing also remains important. The aggregate patterns point to a common service-dominated structure in both countries.





\label{}




%\paragraph{From unstable jobs to unstable spells}
%As mentioned before, once labour markets feature enough atypical employment, as it is the case in Germany and Spain, it makes sense to consider amalgamating these in-between-states with unemployment when looking at the dynamics labour market states. This is the theme of this paper.


%The unit of measurement is a unique individual, observed in his or her first spell in the year 2012. This year is selected because it gives the best time overlap of the two datasets. Therefore tables \ref{tab:summary_stats} and \ref{tab:labour market stats} correspond to a cross-section of the two administrative panels. The characteristics of the two samples look very similar in Table \ref{tab:summary_stats}. Small differences consist in slightly higher fractions of women and college educated (each) in Germany. These differences are unlikely to drive the differences in duration distributions described below. \\



%Marginal Employment refers to the fraction of individuals with an unstable employment spell in their last observation.

\section{Results}\label{sec:results}

Table \ref{tab:results_main}
presents our baseline variance decomposition. 
The top half of the table shows the total variance and
the estimates of its components. The bottom half of the table shows the share of each component in the total variance. 


\begin{table}[htbb!]
\caption{Baseline Results: variance decomposition of unstable spells}
    \centering
\begin{tabular}{lll}
\toprule
 & Germany & Spain \\
Variable &  &  \\
\midrule
$\sigma^{2}_y$ & 2.568 & 1.855 \\
$\sigma^{2}_c$ & 1.645 & 1.645 \\
$\sigma^{2}_n$ & 0.043 & 0.098 \\
$\sigma^{2}_e$ & 0.694 & -0.049 \\
$\sigma^{2}_b$ & 0.186 & 0.162 \\
\midrule
\% constant & 0.641 & 0.833 \\
\% spell number & 0.017 & 0.053 \\
\% duration dependence & 0.270 & -0.027 \\
\% heterogeneity & 0.072 & 0.087 \\
\hline \hline
N & 81,768 & 69,406 \\
Average duration (days) & 360.175 & 912.447 \\
Std. duration (days) & 494.538 & 1047.165 \\
\bottomrule
\end{tabular}
\label{tab:results_main}
\smallskip
\footnotesize \\
Notes: Duration in log days. Sources: SIAB  (2011-2021), MCVL (2011-2021). \end{table}

In both countries the constant component is the largest: it accounts for 65\% of the variance in Germany and 83\% in Spain.
One key difference between the decomposition in Germany and Spain is that, while Germany exhibits sizeable negative duration dependence (positive $\sigma^{2}_e$), Spain shows a small positive duration dependence (negative $\sigma^{2}_e$). That is, in Germany the longer workers have spent in an unstable spell, the less likely they are to exit to stable employment -- their hazard rate is decreasing over duration, regardless of her type.
Heterogeneity accounts for 9\% of the  variance in Spain, while it is 6\% in Germany. Recall that heterogeneity in hazard rates boosts the variance. In Spain, positive duration dependence and heterogeneity counter each other, while in Germany they reinforce each other. \textcolor{mygreen}{this needs more explanation I think.} These baseline results are robust to controlling for outliers, controlling for age and controlling for education (see Appendix \ref{AppendixRobustness}). 

\textcolor{mygreen}{Explain potential sources of duration dependence, conclude that there are good reasons to suspect these sources very by education levels (e.g. importance of networks, importance of skill depreciation). Done in Table A7. Constant dur dep good support for identifying assumption.}


One key question of our baseline results is if the differences found are due to the different forms of atypical contracts in the two countries (i.e., Germany has mini-jobs and Spain does not, while the incidence of temporary jobs is much higher in Spain than it is in  Germany) or if it is cross-country differences. We address this issue by doing the variance decomposition focusing on spells of temporary contracts in both countries. 
For this exercise, we restrict attention to spells where the majority of time is spent in temporary contracts to ensure sufficient sample size.\footnote{Among spells that contain at least some temporary employment, the median share of time spent in temporary jobs is approximately 0.25 in Germany and 0.45 in Spain. We therefore define ``mostly temporary'' spells as those in which the share of days in temporary employment is at least 0.25 and 0.45 in Germany and Spain respectively.} Table \ref{tab:results_mixed} displays the decomposition for temporary jobs in each country.

We note that the log variance is now much more similar between the two countries than in the baseline. 
The striking result is that these contracts behave very similarly in both countries. Duration dependence in Germany switches from negative to positive and in Spain the magnitude of the effect is now sizeable. All in all, the two countries look more similar. 
In both countries, the longer a worker spends in a temporary contract, the more likely they are to exit into a permanent contract.
These results are robust to considering unstable spells that are only made up of temporary contracts, see appendix table \ref{tab:purespells}.


\begin{table}[htbb!]
    \centering
    \caption{Temporary contracts: variance decomposition of unstable spells}
    \begin{tabular}{lll}
\toprule
 & \multicolumn{1}{c|}{Germany} & \multicolumn{1}{c}{Spain} \\[5pt]
Variable & Mostly Temp & Mostly Temp \\
\midrule
$\sigma^{2}_y$ & 1.305 & 1.505 \\
$\sigma^{2}_c$ & 1.645 & 1.645 \\
$\sigma^{2}_n$ & 0.007 & 0.034 \\
$\sigma^{2}_e$ & -0.398 & -0.410 \\
$\sigma^{2}_b$ & 0.051 & 0.237 \\
\midrule
\% constant & 0.650 & 0.548 \\
\% spell number & 0.005 & 0.022 \\
\% duration dependence & -0.305 & -0.273 \\
\% heterogeneity & 0.039 & 0.157 \\
\hline \hline
N & 5,342 & 27,191 \\
Average duration (days) & 588.105 & 850.299 \\
Std. duration (days) & 529.472 & 960.101 \\
\bottomrule
\end{tabular}
\label{tab:results_mixed}
\smallskip
\footnotesize \\
Notes: Duration in log days. Sources: SIAB  (2011-2021), MCVL (2011-2021). \end{table}


Temporary contracts behave similarly in both countries. To have the full picture, we dig into the other components of unstable spells, namely, unemployment in both countries and 
mini-jobs in Germany. Tables \ref{decomp_mixed_DE} and \ref{decomp_mixed_ES} display the decomposition for these. 
Unemployment spells in 
both countries display negative duration dependence, although the magnitudes are different, this being higher in Spain.
Mini-jobs are completely different to temporary contracts in Germany, featuring negative duration depedence. Mini-jobs in Germany contribute to the baseline results being different in the two countries.



The incidence of atypical contracts is higher for females than males so we next implement the decomposition by sex. 
Table \ref{decomp_sex} shows the results. Interestingly, there is no major differences in each country. If anything, the results hint to the fact that for males the unstable state is more persistent.  

   
\begin{table}[htbb!]
    \centering
    \caption{Decomposition by sex}
\begin{tabular}{lllll}
\toprule
 & \multicolumn{2}{c|}{Germany} & \multicolumn{2}{c}{Spain} \\[5pt]
Variable & female & male & female & male \\
\midrule
$\sigma^{2}_y$ & 2.272 & 2.818 & 1.818 & 1.884 \\
$\sigma^{2}_c$ & 1.645 & 1.645 & 1.645 & 1.645 \\
$\sigma^{2}_n$ & 0.018 & 0.077 & 0.086 & 0.108 \\
$\sigma^{2}_e$ & 0.597 & 0.770 & -0.058 & -0.044 \\
$\sigma^{2}_b$ & 0.012 & 0.326 & 0.144 & 0.176 \\
\midrule
\% constant & 0.724 & 0.584 & 0.842 & 0.826 \\
\% spell number & 0.008 & 0.027 & 0.047 & 0.057 \\
\% duration dependence & 0.263 & 0.273 & -0.032 & -0.023 \\
\% heterogeneity & 0.005 & 0.116 & 0.079 & 0.093 \\
\hline \hline
N & 40,691 & 41,077 & 29,774 & 39,632 \\
Average duration (days) & 378.578 & 341.946 & 899.500 & 922.174 \\
Std. duration (days) & 494.494 & 493.910 & 1033.914 & 1056.912 \\
\bottomrule
\end{tabular}
\label{decomp_sex}
\smallskip
\footnotesize \\
Notes: Duration in log days. Sources: SIAB  (2011-2021), MCVL (2011-2021).
\end{table}
    
Our results can be affected by the business cycle. Table \ref{tab:results_spain_BC} displays the decomposition controlling for changes in unemployment rates or GDP.
In Germany, results are similar to the baseline. Business cycle controls explain around 2\% of the variance, relative to 1\% in the baseline. In Spain, 
controlling for changes in unemployment rates 
play an important role, explaining around 8\%, relative to the 5\% in the baseline. This has to do with the pronounced cyclicality of unemployment rates in Spain.
Duration dependence becomes more prominent, as the controls capture
the effects of bad times on exiting the unstable spell.


\begin{table}[htbb!]
    \centering
    \caption{Decomposition with Business cycle controls}
    \resizebox{\linewidth}{!}{
\begin{tabular}{lllllll}
\toprule
 & \multicolumn{3}{c|}{Germany} & \multicolumn{3}{c}{Spain} \\[5pt]
Variable & Baseline & $\Delta$u & $\Delta$GDP & Baseline & $\Delta$u & $\Delta$GDP \\
\midrule
$\sigma^{2}_y$ & 2.568 & 2.568 & 2.568 & 1.855 & 1.855 & 1.855 \\
$\sigma^{2}_c$ & 1.645 & 1.645 & 1.645 & 1.645 & 1.645 & 1.645 \\
$\sigma^{2}_n$ & 0.043 & 0.069 & 0.059 & 0.098 & 0.141 & 0.102 \\
$\sigma^{2}_e$ & 0.694 & 0.660 & 0.676 & -0.049 & -0.076 & -0.051 \\
$\sigma^{2}_b$ & 0.186 & 0.195 & 0.188 & 0.162 & 0.146 & 0.160 \\
\midrule
\% constant & 0.641 & 0.641 & 0.641 & 0.833 & 0.805 & 0.831 \\
\% spell number and BC controls & 0.017 & 0.027 & 0.023 & 0.053 & 0.076 & 0.055 \\
\% duration dependence & 0.270 & 0.257 & 0.263 & -0.027 & -0.041 & -0.028 \\
\% heterogeneity & 0.072 & 0.076 & 0.073 & 0.087 & 0.079 & 0.086 \\
\hline \hline
N & 81,768 & 81,768 & 81,768 & 69,406 & 69,406 & 69,406 \\
Average duration (days) & 360.175 & 360.175 & 360.175 & 912.447 & 912.447 & 912.447 \\
\bottomrule
\end{tabular}
}
\label{tab:results_spain_BC}
\end{table}



\newpage

\section{Interpretation of Duration Dependence Results}

Duration dependence plays an important role in our results. In this section, we illustrate the economic meaning behind it and shed further light on its qualitative importance. To do so, we focus on the expected number of additional months in an unstable spell under different scenarios. 
We map the estimated decomposition moments
for each country 
to the distribution parameters to allow for 3 possible cases. See Appendix \ref{TecnicalAppendix}
for details.

\begin{enumerate}
    \item \emph{Exponential}: An exponential function that corresponds to a constant hazard rate independent of duration.
    \item \emph{Heterogeneity only}: A mixture of two exponentials to isolate the effect of heterogeneity in exit rates, where $\lambda_{\text{fast}}$
and $\lambda_{\text{slow}}$
parameterize the hazard rates for fast and slow exiters, respectively. We assume that individuals only differ in their underlying exit rates. 
%with $k=1$ in the weibull

\item  \emph{Heterogeneity + DD}:  
We augment case (2) by allowing duration dependence with a Weibull distribution in which the duration dependence parameter $k$ will be different from 1 (as $k=1$ corresponds to an exponential, hence no duration dependence). We assume that individuals differ in their underlying exit rates and, moreover, the hazard rates of leaving 
are a function of duration.
More specifically, in Germany the parameter $k$ that matches the data is below $1$ (hazard falls with duration), corresponding to negative duration dependence; and in Spain the $k$ that matches the data is larger than but very close to $1$ (hazard rises slightly with duration). 

\end{enumerate}

Figure \ref{exp_dur_base} displays the expected additional months in an unstable spell by elapsed duration under the three scenarios above.  

The exponential fit shows the difference in levels of the survival rates in an unstable spell in Germany and Spain. In Spain, the expected time in an unstable spell is 2.5 times higher than in Germany. 

To shed light on the economic meaning of duration dependence, we explore the difference in the profiles of \emph{Heterogeneity only} and \emph{Heterogeneity +DD}. 
As expected from Table \ref{tab:results_main}, in Germany, 
the \emph{Heterogeneity + DD} plot is in general above the  \emph{Heterogeneity only} plot, confirming that there is negative duration dependence. 
For instance, after spending 24 months in an unstable spell, the expected remaining duration is 13.39 additional months under \emph{Heterogeneity only}, and 16.76 additional months under \emph{Heterogeneity + DD}.


As expected from Table \ref{tab:results_main}, in Spain, the \emph{Heterogeneity only} plot is very close to the \emph{Heterogeneity + DD} and, if anything, is slightly above, 
confirming that there is a small but positive duration dependence. 
For instance, after spending 24 months in an unstable spell, the expected remaining duration is 
29.99 additional months under \emph{Heterogeneity only}, and  29.63 additional months under \emph{Heterogeneity + DD}.

For an even more in-depth within country differences, appendix 
figure \ref{het_DD_base} zooms into  the difference in expected months in an unstable spell relative to an individual at the beginning of the unstable spell.




\begin{figure}[h!]
    \centering
    \caption{Expected months in an unstable spell}
    \vspace{-0.5cm}
\begin{center}
\includegraphics[width=0.8\textwidth]{imgs/exp_dur_base_a.png}
\end{center}
\label{exp_dur_base}
\end{figure}

Thus, there are two key distinguishing features between Germany and Spain. On the one hand, the overall exit rate of an unstable spell is much higher in Germany. On the other hand, being unstable is a ``trap'': once you are in the unstable spell, you are stuck and the hazard of exiting falls with duration. In Spain, however, it is the opposite: the exit rates out of unstable spells is low, but there is no strong "trap" effect.


%E.g. after 12 months of time already elapsed in Spain it means under Het and DD that you have double the expected time in unstable than a new entrant, in Germany under Het and DD 7 times. However, if we compare that to het only, in Germany after 12 months you are only two times less likely to leave.
%In Spain, the dark bar is below the orange bar, DD has a negative contribution. No strong “trap” effect from DD itself; most of the observed change over duration comes from heterogeneity/selection. In Germany, the dark blue bar is much higher than the light blue bar, showing the positive contribution of DD. The hazard of exiting an unstable spell falls with duration, so workers stuck longer in instability become “trapped”: the longer the spell, the more extra months they expect to remain unstable.
%The gap between dark and light bars for each country quantifies how much of this “trap” is due to true negative duration dependence, not just compositional changes.


Figure \ref{exp_dur_T} repeats the above analysis for spells made up of temporary contracts in both countries. This is in line with the results from Figure \ref{tab:results_mixed}, giving a qualitative interpretation to the meaning behind it. Intriguingly, in both Germany and Spain, the plot of \emph{Heterogeneity + DD} lies below that of \emph{Heterogeneity Only}; there is positive duration dependence counteracting the effect of heterogeneity. \textcolor{mygreen}{It seems counterintuitive to call positive duration dependence a situation in which exit rates decrease with spell duration. } For instance, after spending 24 months in an unstable spell, the expected remaining duration is 17.89 under \emph{Heterogeneity + DD} and an additional 18.85 under \emph{Heterogeneity only}. In Germany, after spending 24 months in an unstable spell, the expected remaining duration is 16.64 under \emph{Heterogeneity + DD} and an additional 17.45 under \emph{Heterogeneity only}. Hence, when focusing on only temporary contracts, we find that there is no strong ``trap'' effect in either country. However, turning to the exponential fit, restricting attention to temporary contracts substantially increases survival rates in an unstable spell in Germany, while it reduces them in Spain, relative to the baseline analysis. Thus, the particular atypical contract under consideration appears to be a key factor in determining the relative importance of duration dependence versus heterogeneity.



\\

\begin{figure}[h!]
    \centering
    \caption{Expected months in an unstable spell, Temp Only}
    \vspace{-0.5cm}
\begin{center}
\includegraphics[width=0.8\textwidth]{imgs/exp_dur_T_a.png}
\end{center}
\label{exp_dur_T}
\end{figure}




\newpage
\clearpage
\section{Sample Selection}

\textcolor{magenta}{Need to cite appendix A.3 or remove it}

% 1. The problem--part 1: why do we have sample selection? we require 2 spells

The requirement of two completed unstable spells imposes a sample restriction: for some individuals, we never observe two completed spells. This could be for two reasons: 
\textcolor{blue}{First, the finite observation window. Within a finite observation window, a worker who cycles quickly between unstable and stable spells is more likely to accumulate two complete cycles than a worker who cycles slowly. Hence, the estimation sample over-represents fast-cycling individuals. Secondly, heterogeneity in cycling behaviour. Some workers are of a low-type and never accumulate unstable spells. These channels interact as it is the low-cycling, more heterogenous, workers who are most likely to be censored. Our calibration exercise separates them.}the restricted time frame or more heterogeneity than we are able to capture.\\

To see this, consider a worker who goes through 2 unstable spells when young, but after that, she gets a stable job that lasts 15 years. If we use our 10-year sample, we could only observe the last stable spell. But if we increase the time window to 20 years, then we can observe the two unstable spells she had prior to, so she comes into our sample. Consider also Figure \ref{spell_selection}, the case of a worker who has two very long, unstable spells of 5 years each. With the 10-year time window (sample 1), we may only be able to observe the latter of these two spells, but after extending it (sample 2), we capture the earlier spell - this corresponds to the case of a \textit{slow} worker that has lower hazards than the rest. This case illustrates how sample selection can affect our results: we are missing the \textit{slow} workers.\\

\begin{figure}[h!]
    \centering
    \caption{Extending the Sample}
    \vspace{-0.7cm}
\begin{center}
\includegraphics[width=0.8\textwidth]{imgs/spell_selection.png}
\end{center}
\label{spell_selection}
\end{figure}


There is a concern that our sample is extremely selected - that is, that the workers we are capturing have very high hazards of leaving unstable spells \textcolor{red}{and/or leaving stable spells. Concern: leave too fast the stable spell}. In other words, they are \textit{fast}, as their spells are on average very short, compared to the larger population that may have longer spells. This is particularly relevant for Spain, as we find some negative duration dependence. Not so much for Germany, where we already find some positive duration dependence - which is consistent with an overall \textit{faster} population.\\

To see if that is the case, we consider two approaches: increasing time windows (so we have more chances of capturing \textit{slow} workers) and using the information we have on all workers to infer a reasonable underlying hazard distribution.\\

The first approach is straightforward. For Spain, we increase the time window allowing spells to start from 2005, 1999, and 1992.\footnote{We also did 1996, but the results are very similar to 1992, which is the year that temporary contracts start being recorded.} Increasing the time window means that we have more time to capture transitions in and out of unstable spells. We find that the results from increased time windows are not dramatically different from our main decomposition (\textcolor{magenta}{see appendix XXXX with the results}), and while heterogeneity and duration dependence increase in relevance, the constant component is still the most prominent and duration dependence is still close to zero. These observations are consistent with the view that we are not missing ``true heterogeneity", but that the missing individuals are very similar to the ones in our sample. When unstable spells last long on average (like in Spain), it takes a longer time window to capture two completed cycles -- the individuals are \textit{randomly} missing. \textcolor{blue}{In Germany, we cannot extend the time window beyond 2011 as temporary contracts are only recorded in the data from then on. Instead, we repeat the decomposition for shorter time windows for a comparison of the results. Table A.13 in appendix (hyperlink here). Results are not dramatically different etc.}\\

\textcolor{blue}{This is n o longer up to date. See sample selection note instead.}

The second approach is more nuanced. Suppose that instead of slow and fast individuals, we have a \textit{continuum} of hazard rates among the population. That is, each individual (as our method assumes) has an intrinsic hazard rate. Then there will be a distribution of hazards among the population. These are hazards, so someone with a very low hazard of finishing an unstable spell (slow worker) would be impossible to capture, even if we extended the time window for a very long time -- in expectation. Some of these very slow workers will be lucky enough to leave once, never to return again.\\

The idea is then to use all of the data, not only our sample, to infer what reasonable underlying distribution could result in our observed sample of 2 completed spells, 1 completed and one incomplete spell only. That is, if most of our sample had a very low hazard, then it would be very unlikely that we would get a 2 completed spell sample. Or in other words, the likelihood of observing a large share of workers with 2 completed spells would be close to zero. The more workers we have with a high hazard rate, the larger the share we expect to see with 2 completed spells.\\

We then proceed as follows:
\begin{enumerate}
    \item Find the shares in the data with 0 (only one incomplete spell), 1 (one completed but not two) and 2 completed spells.
    \item Use a simple duration model to recover the underlying average hazard - including incomplete spells. \textcolor{red}{ delta or lamba?}
    \item Using this average, build a distribution of underlying hazards targeting the same average but letting the variance vary - and/or shape parameters, depending on the distribution choice.
    \item Use this estimated distribution to simulate a synthetic panel of spells, and measure the shares in this panel with 0, 1 and 2 completed spells.
    \item Plot the resulting hazard distributions for each group (the ones with 0, 1 and 2 completed spells) and compare how different they are. This will give us an estimate of how close is the distribution of hazards in our sample (2 completed spells) with the underlying distribution
\end{enumerate}

The most involved step is step 3. For this, we need to select a random variable distribution family, and conditional on it, change the parameters (like the mean and variance) to match the mean hazard and the share of workers with 0, 1 and 2 spells. For the selection of the underlying distribution, we consider positive continuous distributions, as hazard rates are positive. \\

\textcolor{magenta}{[Insert here justification for normal - lognormal transformation and mention Shimer and Borovickova time stopping paper]}\\



% 2. The problem--part 2: This is aggravated by a restricted time-window of our sample

% 3. We want to see how this affects our results.

% 4. What do we do?
% Show the  mapping between the underlying distribution of job finding rates (of the population) and the distribution in the selected sample.

% Based on the transitions that we observe (probability leave Uns and stable), we can infer how likely we are to miss those without 2 spells.





% 5. Details of how we do (4)

% 6. How we solve the problem: a measure of how our selection problem is biasing our results.

% 7. Descriptive stats for sample selection: add them?

% 8. Extra: Decomposition with: Take the 30 year sample, and select workers in the next 10 years only. Separate longer window vs. different period.

Figure \ref{fig:hazard_sim10} shows the result of this procedure. In black we have the estimated underlying distribution of hazards among the population. This represents the ``true type'' in the population. The green distrbituion shows the distribution of our selected sample - that is, those who complete two spells. This distribution captures some \textit{slow} workers (the ones with hazards close to zero) but fewer workers than in the underlying population. At higher hazards, the opposite is true: the two-spell distribution has more density than the underlying distribution. Fast types are over-represented. In the 10-year sample, the 1-sample spell is arguably a better representation than the whole population - but we cannot capture it in our decomposition because the two-spell requirement and the 10-year windown.\\


\begin{figure}[h!]
\centering
 \includegraphics{selection/selection_sim_10year.png}
 \caption{Distribution hazard completing one unstable spell 10 year horizon}
 \label{fig:hazard_sim10}
\end{figure}    

how different 1 spell vs 2 spell.\\
how respresentative is 2 spell wrt to pop

here with both, cristina redo 
\begin{figure}[h!]
\centering
 \includegraphics{selection/selection_sim_both_10years.png}
 \caption{Distribution hazard completing one unstable spell 10 year horizon}
 \label{fig:hazard_sim10}
\end{figure} 

In contrast, Figure \ref{fig:hazard_sim20} shows that with a longer timeframe, the 2-spell distribution looks much closer to the data than the 1 or 0 spell distributions. This illsutrates how increasing the time window (which we can do in Spain) allows us to include shares of \textit{slow} workers that look much more like the data.\\

But as noted before, increasing the time window comes with its drawbacks. First, we cannot really compare with Germany anymore, as in Germany, we only have a 10-year window. Secondly, the further back into the past we go with the Spanish data, the less representative of the population at that time (say 1999) it becomes, as we are observing the sample of 2020 individuals as they were in their past. So we observe the attached workers who "survive" in the labour market until 2020. We also lose some data quality beyond 1999, so we recommend some caution interpreting results beyond this date. 

\begin{figure}[h!]
\centering
\includegraphics{selection/selection_sim_20year.png}
    \caption{Distribution hazard completing one unstable spell 20 year horizon}
\label{fig:hazard_sim20}
\end{figure}    

% Spain, time windows: basic results now in appendix.

       
\newpage
\section{Conclusion}\label{sec:discussion}


Our analysis reveals that, despite Germany and Spain exhibiting similar survival rates of unstable employment, the underlying mechanisms driving these dynamics differ significantly. In Germany, duration dependence is negative, meaning workers in unstable jobs become more likely to transition to stable employment over time. In contrast, Spain exhibits positive duration dependence, where the likelihood of leaving instability decreases the longer a worker remains in an unstable job.  These differences are closely linked to the structure of atypical employment in each country—Germany’s mini-jobs, which offer low pay but continuous employment, versus Spain’s temporary contracts, which provide higher wages but strict time limits and frequent terminations.  
Our findings challenge the idea that stronger attachment to atypical employment leads to greater employment stability and emphasize the role of labor market institutions in shaping long-term employment outcomes.  

Additionally, while our methodology does not impose strong functional form assumptions, it relies on observing at least two completed unstable spells per individual, which may introduce sample selection concerns. The use of comprehensive administrative data helps mitigate these issues, and we show that increasing the observation period (larger T) reduces potential biases. 




\clearpage
\begin{spacing}{0.8}
\currentpdfbookmark{References}{Refbookmark}
\bibliographystyle{chicago}
\bibliography{biblio}
\end{spacing}

\newpage
\begin{appendices}
\renewcommand{\thetable}{A.\arabic{table}}
\renewcommand{\thesection}{\Alph{section}}
\renewcommand{\thefigure}{A.\arabic{figure}}
\setcounter{table}{0}
\setcounter{figure}{0}

\section{Technical Appendix }\label{TecnicalAppendix}

\subsection{Simulations}
\label{app_simulations}
\textcolor{magenta}{Include: Simulated survival function}\\ 

This simulation illustrates the difficulty in identifying separately DD from heterogeneity (large literature) but that this method can naturally do.

\subsection{Derivation of the Baseline Variance} 
\label{baseline_var}






This appendix derives the result that the variance of log durations under an exponential distribution equals $\pi^2/6 \approx 1.645$, independent of the rate parameter. This constant appears as $\sigma^2_c$ in our variance decomposition and represents the irreducible variance that remains even in the absence of unobserved heterogeneity or duration dependence. Appendix \ref{weibull} builds on this result to derive the closed-form solution for the Weibull mixture model. \\ 

Let $T \sim \text{Exponential}(\lambda)$, with density $f(t) = \lambda e^{-\lambda t}$ for $t > 0$. Define $X = \log T$. We derive $\text{Var}(X)$.


We have
\begin{equation}
    \mathbb{E}[\log T] = \int_0^\infty \log(t)\, \lambda e^{-\lambda t}\, dt.
\end{equation}
Substituting $u = \lambda t$ yields
\begin{equation}
    \mathbb{E}[\log T] = \int_0^\infty \log\!\left(\frac{u}{\lambda}\right) e^{-u}\, du = \int_0^\infty \log(u)\, e^{-u}\, du - \log\lambda.
\end{equation}
The integral $\int_0^\infty \log(u)\, e^{-u}\, du = \Gamma'(1)$ is the derivative of the Gamma function evaluated at 1. Using the digamma function $\psi(s) = \Gamma'(s)/\Gamma(s)$ and the classical result $\psi(1) = -\gamma$ \citep{abramowitz1965}, where $\gamma \approx 0.5772$ is the Euler--Mascheroni constant, we obtain $\Gamma'(1) = -\gamma$. Therefore:
\begin{equation}\label{eq:mean_logT}
    \mathbb{E}[\log T] = -\gamma - \log\lambda.
\end{equation}


Applying the same substitution:
\begin{align}
    \mathbb{E}[(\log T)^2] &= \int_0^\infty (\log u - \log\lambda)^2\, e^{-u}\, du \notag \\
    &= \underbrace{\int_0^\infty \log^2(u)\, e^{-u}\, du}_{\Gamma''(1)} - 2\log\lambda \underbrace{\int_0^\infty \log(u)\, e^{-u}\, du}_{-\gamma} + \log^2\lambda \underbrace{\int_0^\infty e^{-u}\, du}_{1} \notag \\
    &= \Gamma''(1) + 2\gamma\log\lambda + \log^2\lambda, \label{eq:second_moment}
\end{align}
where $\Gamma''(1)$ follows from differentiating the Gamma function twice \citep{abramowitz1965}.



Computing $\text{Var}(\log T) = \mathbb{E}[(\log T)^2] - (\mathbb{E}[\log T])^2$:
\begin{align}
    \text{Var}(\log T) &= \Gamma''(1) + 2\gamma\log\lambda + \log^2\lambda - (\gamma + \log\lambda)^2 \notag \\
    &= \Gamma''(1) + 2\gamma\log\lambda + \log^2\lambda - \gamma^2 - 2\gamma\log\lambda - \log^2\lambda \notag \\
    &= \Gamma''(1) - \gamma^2. \label{eq:var_intermediate}
\end{align}
All terms involving $\lambda$ cancel exactly. The variance of $\log T$ is a pure property of the exponential distributional shape, independent of the rate parameter.



The second derivative of the Gamma function at 1 can be expressed using the polygamma functions \citep{abramowitz1965}:
\begin{equation}
    \Gamma''(1) = \left[\psi'(1) + \psi(1)^2\right] \Gamma(1) = \psi'(1) + \gamma^2,
\end{equation}
where $\psi'(s) = \sum_{n=0}^{\infty} (n+s)^{-2}$ is the trigamma function. Evaluating at $s = 1$:
\begin{equation}\label{eq:basel}
    \psi'(1) = \sum_{n=1}^{\infty} \frac{1}{n^2} = \frac{\pi^2}{6},
\end{equation}
which is the solution to the Basel problem \citep{euler1740}. Substituting into \eqref{eq:var_intermediate}:
\begin{equation}\label{eq:var_logT_final}
   \;\text{Var}(\log T) = \Gamma''(1) - \gamma^2 = \frac{\pi^2}{6} + \gamma^2 - \gamma^2 = \frac{\pi^2}{6} \approx 1.6449.\;
\end{equation}



For a $K$-type mixture of exponentials with type-specific rates $\{\lambda_i\}_{i=1}^K$ and mixing probabilities $\{p_i\}_{i=1}^K$, the law of total variance yields:
\begin{equation}\label{eq:decomp_exponential}
    \sigma^2_y = \underbrace{\frac{\pi^2}{6}}_{\sigma^2_c} + \underbrace{\sum_{i=1}^K p_i \left(\log\frac{1}{\lambda_i} - \sum_{j=1}^K p_j \log\frac{1}{\lambda_j}\right)^2}_{\sigma^2_b}.
\end{equation}

This is a direct application of the variance decomposition 
in equation~\eqref{eq:main_decomp} in the main text. The within-type variance 
$\sigma^2_w = \frac{1}{N}\sum_{i=1}^N \sigma^2_i$ reduces to 
$\pi^2/6$ because each exponential type has identical log-duration 
variance regardless of its rate parameter $\lambda_i$, as shown 
above. The between-type variance $\sigma^2_b = \frac{1}{N}\sum_{i=1}^N 
(\mu_i - \mu_y)^2$ becomes $\text{Var}_i[\log(1/\lambda_i)]$ 
since the type-specific means of $\log T$ differ only through 
$\log(1/\lambda_i)$. The connection to extreme value theory is direct: $\log T$ follows a Gumbel distribution with unit scale when $T$ is exponential, and a Gumbel random variable with scale 1 has variance $\pi^2/6$ 

The within-type component $\sigma^2_c = \pi^2/6$ is irreducible: it is the variance that remains even when all types are identical, i.e., when there is no heterogeneity. Any excess variance $\sigma^2_y > \pi^2/6$ must originate from between-type heterogeneity ($\sigma^2_b > 0$) or duration dependence ($k \neq 1$, which shifts the within-type variance away from $\pi^2/6$). Our decomposition in the main text separates exactly these contributions.
When we generalise to the Weibull distribution in Appendix \ref{weibull}, the shape parameter $k$ rescales this baseline variance to $\pi^2/(6k^2)$, reflecting the role of duration dependence. 
\citep{embrechts1997}. 





\subsection{Analytic Solution for Weibull Mixture Model}
\label{weibull}
Here we derive the closed-form solution for a same-shape Weibull mixture with two unobserved types at equal proportions. The structural parameters can be recovered directly from population moments without numerical optimisation.


Suppose durations are generated by a two-type mixture of Weibull distributions with a common shape parameter $k$ but heterogeneous scale parameters. Specifically, let the two types have scale parameters $\lambda_0$ and $\lambda_0 \alpha_0$, i.e. $\lambda_{\text{slow}}$ and $\lambda_{\text{fast}}$, with $\alpha > 1$,  each occurring with probability $\frac{1}{2}$. As derived in Appendix \ref{baseline_var} if $T \sim \text{Weibull}(k, \lambda)$, then $\log T$ follows a Gumbel (Type I extreme value) distribution with location $\log \lambda$ and scale $1/k$. Its moments are:
\begin{align}
    \mathbb{E}[\log T] &= \log \lambda + \frac{\gamma}{k}, \label{eq:gumbel_mean} \\
    \text{Var}(\log T) &= \frac{\pi^2}{6k^2}, \label{eq:gumbel_var}
\end{align}
where $\gamma \approx 0.5772$ is the Euler--Mascheroni constant (\cite{johnson1994continuous}).
The four moments of the mixture map to the structural parameters as shown in Table~\ref{tab:moments}.

\begin{table}[htbp!]
\centering
\caption{Moment--Parameter Mapping for the Weibull Mixture}
\label{tab:moments}
\begin{tabular}{@{}lll@{}}
\toprule
Moment & Formula & Interpretation \\
\midrule
$\sigma^2_{w}$ & $\dfrac{\pi^2}{6k^2}$ & Within-type (duration dependence) \\[10pt]
$\sigma^2_{b}$ & $\dfrac{1}{4}\left(\log \alpha_0\right)^2$ & Between-type (unobserved heterogeneity) \\[10pt]
$\sigma^2_{e}$ & $\dfrac{\pi^2}{6}\!\left(\dfrac{1}{k^2} - 1\right)$ & Residual variance \\[10pt]
$\mathbb{E}[T]$ & $\lambda_0 \cdot \Gamma\!\left(1 + \dfrac{1}{k}\right) \cdot \dfrac{1 + \alpha_0}{2}$ & Mean duration \\[6pt]
\bottomrule
\end{tabular}
\end{table}



The three structural parameters $(k, \alpha_0, \lambda_0)$ can be recovered directly from the data moments $(\mu_y, \sigma^2_{b}, \sigma^2_{e})$ as follows:
\begin{align}
    k &= \frac{\pi}{\sqrt{6 \cdot \sigma^2_{w}}}, \qquad \text{where } \sigma^2_{w} = \sigma^2_{e} + \frac{\pi^2}{6}, \label{eq:k_solution} \\[8pt]
    \alpha_0 &= \exp\!\left(2\sqrt{\sigma^2_{b}}\right), \label{eq:alpha_solution} \\[8pt]
    \lambda_0 &= \frac{2\,\mu_y}{\Gamma\!\left(1 + \frac{1}{k}\right)\left(1 + \alpha_0\right)}. \label{eq:lambda_solution}
\end{align}


An important feature of this system is that it is over-identified. There are four target moments: $\mu_y$, $\sigma^2_{y}$, $\sigma^2_{b}$, and $\sigma^2_{e}$, but only three free parameters. The model imposes the restriction:
\begin{equation}
    \sigma^2_{y} = \sigma^2_{w} + \sigma^2_{b},
\end{equation}
which is a direct consequence of the law of total variance (\cite{heckman1984method}). All four moments can be matched simultaneously only if the data conform exactly to the Weibull mixture. The analytic solution presented here uses $(\mu_y, \sigma^2_{b}, \sigma^2_{e})$ to solve exactly for the structural parameters via equations \eqref{eq:k_solution}--\eqref{eq:lambda_solution}, and then lets $\sigma^2_{y}$ be whatever the model implies. This identifies the structural parameters governing heterogeneity ($\sigma^2_{b}$), duration dependence ($\sigma^2_{e}$), and the level of durations ($\mu_y$), while accepting that total variance may not be perfectly matched if the data-generating process departs from the assumed Weibull mixture. The over-identifying restriction $\sigma^2_{y} = \sigma^2_{w} + \sigma^2_{b}$ then serves as a specification test: a large discrepancy between the model-implied and empirical total variance indicates misspecification of the distributional assumption \cite{van2001duration, abbring2007econometric}.

\clearpage

\section{Appendix: Descriptive Statistics by Industry}\label{AppendixIndustry}



\begin{table}[httb!]
    \centering
    \scriptsize
    \caption{Industry of last employment, Germany}
\begin{tabular}{lrr}
\toprule
Industry & $N$ & Share (\%) \\
\midrule
Missing & 7,500 & 9.17 \\
Agriculture & 1,663 & 2.03 \\
Manufacturing (consumer goods) & 3,309 & 4.05 \\
3 & 3,811 & 4.66 \\
Manufacturing (machinery \& electronics) & 4,622 & 5.65 \\
Construction & 5,019 & 6.14 \\
Wholesale and Retail Trade & 9,453 & 11.56 \\
Transportation and Storage & 4,888 & 5.98 \\
Hospitality & 3,698 & 4.52 \\
Business Services & 17,126 & 20.94 \\
Financial and Insurance & 2,092 & 2.56 \\
Public Admin.\ and Defence & 3,809 & 4.66 \\
Education & 2,322 & 2.84 \\
Health Services & 12,456 & 15.23 \\
\midrule
Total & 81,768 & 100.00 \\
\bottomrule
\end{tabular}

\label{ind_DE}
\vspace{0.5cm}

    \caption{Industry of last employment, Spain}
\begin{tabular}{lrr}
\toprule
Industry & $N$ & Share (\%) \\
\midrule
Missing & 538 & 0.78 \\
Agriculture & 1,167 & 1.68 \\
Mining and Quarrying & 73 & 0.11 \\
Manufacturing (primary) & 2,689 & 3.87 \\
Manufacturing (machinery) & 3,450 & 4.97 \\
Manufacturing (detail and repair) & 1,099 & 1.58 \\
Energy, Gas and Waste Treatment & 537 & 0.77 \\
Construction & 6,610 & 9.52 \\
Retail and Repairs & 13,366 & 19.26 \\
Transport and Storage & 5,426 & 7.82 \\
Hospitality & 8,930 & 12.87 \\
Communication and Programming & 2,860 & 4.12 \\
Financial Activities & 935 & 1.35 \\
Real Estate & 780 & 1.12 \\
Professional and Scientific Services & 4,290 & 6.18 \\
Auxiliary Services (cleaning, security) & 5,901 & 8.50 \\
Public Administration & 861 & 1.24 \\
Education & 1,599 & 2.30 \\
Health and Social Services & 3,429 & 4.94 \\
Other Services & 4,866 & 7.01 \\
\midrule
Total & 69,406 & 100.00 \\
\bottomrule
\end{tabular}
\label{ind_ES}
\end{table}


\clearpage

\section{Appendix: Robustness checks}\label{AppendixRobustness}




\subsection{Appendix: Decomposition, No Outliers}
Note: outliers are people with spells 10 years or longer or more than 12 jobs with different employers per spell.

\begin{table}[httb!]
    \centering
    \caption{Main decomposition, no outliers}
\begin{tabular}{lll}
\toprule
 & Germany & Spain \\
Variable &  &  \\
\midrule
$\sigma^{2}_y$ & 2.563 & 1.792 \\
$\sigma^{2}_c$ & 1.645 & 1.645 \\
$\sigma^{2}_n$ & 0.042 & 0.074 \\
$\sigma^{2}_e$ & 0.690 & -0.109 \\
$\sigma^{2}_b$ & 0.186 & 0.182 \\
\midrule
\% constant & 0.642 & 0.796 \\
\% spell number & 0.016 & 0.041 \\
\% duration dependence & 0.269 & -0.061 \\
\% heterogeneity & 0.072 & 0.102 \\
\hline \hline
N & 81,656 & 64,107 \\
Average duration (days) & 357.992 & 803.791 \\
Std. duration (days) & 485.860 & 888.478 \\
\bottomrule
\end{tabular}

\end{table}

\clearpage

\subsection{Appendix: Decomposition by age}


\begin{table}[htbb!]
    \centering
    \caption{Decomposition by age group}
   \begin{tabular}{lllll}
\toprule
 & \multicolumn{2}{c|}{Germany} & \multicolumn{2}{c}{Spain} \\[5pt]
Variable & 25-35 & 35-55 & 25-35 & 35-55 \\
\midrule
$\sigma^{2}_y$ & 2.514 & 2.497 & 1.657 & 1.797 \\
$\sigma^{2}_c$ & 1.645 & 1.645 & 1.645 & 1.645 \\
$\sigma^{2}_n$ & 0.046 & 0.017 & 0.003 & 0.075 \\
$\sigma^{2}_e$ & 0.757 & 0.469 & -0.191 & -0.078 \\
$\sigma^{2}_b$ & 0.066 & 0.366 & 0.200 & 0.155 \\
\midrule
\% constant & 0.654 & 0.659 & 0.762 & 0.829 \\
\% spell number & 0.018 & 0.007 & 0.002 & 0.042 \\
\% duration dependence & 0.301 & 0.188 & -0.115 & -0.043 \\
\% heterogeneity & 0.026 & 0.146 & 0.121 & 0.086 \\
\hline \hline
N & 54,986 & 19,730 & 23,866 & 27,643 \\
Average duration (days) & 381.364 & 271.504 & 631.152 & 838.058 \\
Std. duration (days) & 488.420 & 437.776 & 610.922 & 863.063 \\
\bottomrule
\end{tabular}
\end{table}

\begin{table}[htbb!]
    \centering
    \caption{Decomposition with age controls}
    \resizebox{0.8\linewidth}{!}{
   \begin{tabular}{lllll}
\toprule
 & \multicolumn{2}{c|}{Germany} & \multicolumn{2}{c}{Spain} \\[5pt]
Variable & Baseline & Age controls & Baseline & Age controls \\
\midrule
$\sigma^{2}_y$ & 2.568 & 2.568 & 1.855 & 1.855 \\
$\sigma^{2}_c$ & 1.645 & 1.645 & 1.645 & 1.645 \\
$\sigma^{2}_n$ & 0.043 & 0.086 & 0.098 & 0.105 \\
$\sigma^{2}_e$ & 0.694 & 0.676 & -0.049 & -0.064 \\
$\sigma^{2}_b$ & 0.186 & 0.161 & 0.162 & 0.169 \\
\midrule
\% constant & 0.641 & 0.641 & 0.833 & 0.818 \\
\% spell number & 0.017 & 0.033 & 0.053 & 0.057 \\
\% duration dependence & 0.270 & 0.263 & -0.027 & -0.034 \\
\% heterogeneity & 0.072 & 0.063 & 0.087 & 0.091 \\
\hline \hline
N & 81,768 & 81,768 & 69,406 & 69,406 \\
Average duration (days) & 360.175 & 360.175 & 912.447 & 912.447 \\
\bottomrule
\end{tabular}
}
\end{table}

Note: Age controls are dummies for age younger or equal to 30, age between 30 and 45, age 45+.




\clearpage


\subsection{Appendix: Decomposition by education}
\begin{table}[htbb!]
    \centering
    \caption{Germany, by educational level}
\begin{tabular}{lllll}
\toprule
 & college & higher & secondary & primary \\
Variable &  &  &  &  \\
\midrule
$\sigma^{2}_y$ & 2.508 & 2.554 & 2.516 & 2.914 \\
$\sigma^{2}_c$ & 1.645 & 1.645 & 1.645 & 1.645 \\
$\sigma^{2}_n$ & 0.152 & 0.022 & 0.022 & 0.167 \\
$\sigma^{2}_e$ & 0.667 & 0.699 & 0.600 & 0.856 \\
$\sigma^{2}_b$ & 0.044 & 0.188 & 0.249 & 0.246 \\
\midrule
\% constant & 0.656 & 0.644 & 0.654 & 0.565 \\
\% spell number & 0.061 & 0.009 & 0.009 & 0.057 \\
\% duration dependence & 0.266 & 0.274 & 0.238 & 0.294 \\
\% heterogeneity & 0.018 & 0.074 & 0.099 & 0.085 \\
\hline \hline
N & 16,798 & 43,785 & 16,437 & 1,749 \\
Average duration (days) & 429.743 & 349.011 & 325.666 & 435.132 \\
Std. duration (days) & 543.028 & 477.321 & 474.211 & 605.592 \\
\bottomrule
\end{tabular}
\end{table}

\begin{table}[htbb!]
    \centering
    \caption{Spain, by educational level}
    \begin{tabular}{lllll}
\toprule
 & college & higher & secondary & primary \\
Variable &  &  &  &  \\
\midrule
$\sigma^{2}_y$ & 1.850 & 1.764 & 1.839 & 1.976 \\
$\sigma^{2}_c$ & 1.645 & 1.645 & 1.645 & 1.645 \\
$\sigma^{2}_n$ & 0.054 & 0.100 & 0.123 & 0.110 \\
$\sigma^{2}_e$ & -0.028 & -0.114 & -0.052 & 0.040 \\
$\sigma^{2}_b$ & 0.179 & 0.133 & 0.123 & 0.181 \\
\midrule
\% constant & 0.859 & 0.804 & 0.838 & 0.833 \\
\% spell number & 0.029 & 0.057 & 0.067 & 0.056 \\
\% duration dependence & -0.015 & -0.064 & -0.028 & 0.020 \\
\% heterogeneity & 0.097 & 0.075 & 0.067 & 0.092 \\
\hline \hline
N & 14,908 & 17,766 & 26,847 & 8,240 \\
Average duration (days) & 723.076 & 880.017 & 999.712 & 1043.674 \\
Std. duration (days) & 848.861 & 972.619 & 1116.282 & 1231.377 \\
\bottomrule
\end{tabular}

\end{table}

Notes: ``Higher" means \textit{bachillerato}, secondary means \textit{secundaria}, primary means \textit{primaria}.

\clearpage

\subsection{Contracts vs countries}

\begin{table}[htbb!]
    \centering
    \caption{Pure Temporary Spells}
    \resizebox{\linewidth}{!}{
    \begin{tabular}{lllll}
\toprule
 & \multicolumn{2}{c|}{Germany} & \multicolumn{2}{c}{Spain} \\[5pt]
Variable & Baseline & Only Temporary jobs & Baseline & Only Temporary Jobs \\
\midrule
$\sigma^{2}_y$ & 2.568 & 1.716 & 1.855 & 1.783 \\
$\sigma^{2}_c$ & 1.645 & 1.645 & 1.645 & 1.645 \\
$\sigma^{2}_n$ & 0.043 & 0.020 & 0.098 & -0.000 \\
$\sigma^{2}_e$ & 0.694 & -0.041 & -0.049 & -0.182 \\
$\sigma^{2}_b$ & 0.186 & 0.092 & 0.162 & 0.320 \\
\midrule
\% constant & 0.641 & 0.911 & 0.833 & 0.718 \\
\% spell number & 0.017 & 0.011 & 0.053 & -0.000 \\
\% duration dependence & 0.270 & -0.024 & -0.027 & -0.102 \\
\% heterogeneity & 0.072 & 0.054 & 0.087 & 0.180 \\
\hline \hline
N & 81,768 & 3,633 & 69,406 & 6,851 \\
Average duration (days) & 360.175 & 524.052 & 912.447 & 391.880 \\
Std. duration (days) & 494.538 & 509.124 & 1047.165 & 473.679 \\
\bottomrule
\end{tabular}

}
\label{tab:purespells}
\end{table}


\begin{table}[htbb!]
    \centering
    \caption{Mixed Spells: Germany}
    \resizebox{\linewidth}{!}{
   \begin{tabular}{lllll}
\toprule
 & \multicolumn{4}{c}{Germany} \\[5pt]
Variable & Baseline & Mostly Mini & Mostly Temporary & Mostly Unemployment \\
\midrule
$\sigma^{2}_y$ & 2.568 & 2.244 & 1.305 & 2.391 \\
$\sigma^{2}_c$ & 1.645 & 1.645 & 1.645 & 1.645 \\
$\sigma^{2}_n$ & 0.043 & 0.035 & 0.007 & 0.002 \\
$\sigma^{2}_e$ & 0.694 & 0.495 & -0.398 & 0.455 \\
$\sigma^{2}_b$ & 0.186 & 0.070 & 0.051 & 0.290 \\
\midrule
\% constant & 0.641 & 0.733 & 0.650 & 0.688 \\
\% spell number & 0.017 & 0.016 & 0.005 & 0.001 \\
\% duration dependence & 0.270 & 0.220 & -0.305 & 0.190 \\
\% heterogeneity & 0.072 & 0.031 & 0.039 & 0.121 \\
\hline \hline
N & 81,768 & 4,725 & 5,342 & 40,207 \\
Average duration (days) & 360.175 & 486.033 & 588.105 & 256.146 \\
Std. duration (days) & 494.538 & 547.622 & 529.472 & 439.904 \\
\bottomrule
\end{tabular}
}
\label{decomp_mixed_DE}
\end{table}


\begin{table}[htbb!]
    \centering
    \caption{Mixed Spells: Spain}
    \resizebox{\linewidth}{!}{
\begin{tabular}{llll}
\toprule
 & \multicolumn{3}{c}{Spain} \\[5pt]
Variable & Baseline & Mostly Temporary & Mostly Unemployment \\
\midrule
$\sigma^{2}_y$ & 1.855 & 1.505 & 2.247 \\
$\sigma^{2}_c$ & 1.645 & 1.645 & 1.645 \\
$\sigma^{2}_n$ & 0.098 & 0.034 & 0.131 \\
$\sigma^{2}_e$ & -0.049 & -0.410 & 0.130 \\
$\sigma^{2}_b$ & 0.162 & 0.237 & 0.341 \\
\midrule
\% constant & 0.833 & 0.548 & 0.732 \\
\% spell number & 0.053 & 0.022 & 0.058 \\
\% duration dependence & -0.027 & -0.273 & 0.058 \\
\% heterogeneity & 0.087 & 0.157 & 0.152 \\
\hline \hline
N & 69,406 & 27,191 & 20,648 \\
Average duration (days) & 912.447 & 850.299 & 823.910 \\
Std. duration (days) & 1047.165 & 960.101 & 1016.173 \\
\bottomrule
\end{tabular}
}
\label{decomp_mixed_ES}
\end{table}



\clearpage


\section{Appendix: Interpretation}

\subsection{Appendix: Baseline}
\begin{figure}[h!]
    \centering
    \caption{Difference in expected  months in an unstable spell}
 
\begin{center}
\centering
   \vspace{-0.5cm}
\makebox[\textwidth]{\includegraphics[width=1.2\textwidth]{imgs/het_DD_base_a.png}
}
\footnotesize \\
\centering
Notes: Relative to an individual at the beginning of the unstable spell.
\end{center}
\label{het_DD_base}
\end{figure}

%Interpretation of figure: For someone who has already been unstable for e.g. 24 months (x-axis is elapsed duration in unstable spell), how does their expected remaining time in unstable compare to a new entrant (under different k specs), and the difference between the two colours at each duration tells us how much of that is due to true duration DD versus heterogeneity in exit rates. 
%The “Heterogeneity only” bars increase with duration, even without true DD, because over time, the high-hazard (i.e., $\lambda_{\text{fast}}$) types leave earlier.
%The remaining pool is increasingly composed of low-hazard ($\lambda_{\text{slow}}$) types. So, the average hazard falls and expected remaining time rises purely from heterogeneity, not individual hazard changing with time. The extra increase in the heterogeneity + DD bars (navy/dark red) compared to heterogeneity-only (orange/light blue) reflects true duration dependence on top of this selection effect.


 
\clearpage

\section{Expanding the time window}
These results are sensitive to our sample period and the inclusion of outliers. Tables \ref{tab:results_spain_sample} and \ref{tab:results_ger_sample} show the trade-off between including more years in the sample, which grants us more observations, but it also increases the heterogeneity of the sample. For Spain, the inclusion of outliers, particularly in the 1999 sample that captures the Great Recession in all of its extent, also boosts the variance and increases the weight of duration dependence.

\begin{table}[htbb!]\label{tab:results_spain_sample}
    \centering
    \caption{Spain, different time windows}
    \begin{tabular}{lcccc}
\toprule
Variable & Baseline & 2005- & 1999- & 1992- \\
\midrule
$\sigma^{2}_y$ & 1.855 & 1.880 & 1.872 & 2.081 \\
$\sigma^{2}_c$ & 1.645 & 1.645 & 1.645 & 1.645 \\
$\sigma^{2}_n$ & 0.098 & 0.009 & 0.008 & 0.001 \\
$\sigma^{2}_e$ & -0.049 & 0.016 & -0.002 & 0.103 \\
$\sigma^{2}_b$ & 0.162 & 0.211 & 0.221 & 0.332 \\
\midrule
\% constant & 0.833 & 0.875 & 0.877 & 0.791 \\
\% spell number & 0.053 & 0.005 & 0.004 & 0.000 \\
\% duration dependence & -0.027 & 0.009 & -0.001 & 0.050 \\
\% heterogeneity & 0.087 & 0.112 & 0.118 & 0.159 \\
\hline \hline
N & 69,406 & 126,285 & 165,237 & 215,907 \\
Average duration (days) & 912.447 & 1023.842 & 1063.454 & 927.491 \\
Std. duration (days) & 1047.165 & 1167.795 & 1195.188 & 1187.574 \\
\bottomrule
\end{tabular}

\end{table}

% \section{Germany: Time Windows}

\begin{table}[htbb!]\label{tab:results_ger_sample}
    \centering
    \caption{Germany, different time windows}
\begin{tabular}{llll}
\toprule
Variable & Baseline & 2014- & 2016- \\
\midrule
$\sigma^{2}_y$ & 2.568 & 2.471 & 2.363 \\
$\sigma^{2}_c$ & 1.645 & 1.645 & 1.645 \\
$\sigma^{2}_n$ & 0.043 & 0.022 & 0.015 \\
$\sigma^{2}_e$ & 0.694 & 0.643 & 0.552 \\
$\sigma^{2}_b$ & 0.186 & 0.161 & 0.152 \\
\midrule
\% constant & 0.641 & 0.666 & 0.696 \\
\% spell number & 0.017 & 0.009 & 0.006 \\
\% duration dependence & 0.270 & 0.260 & 0.234 \\
\% heterogeneity & 0.072 & 0.065 & 0.064 \\
\hline \hline
N & 81,768 & 48,061 & 27,601 \\
Average duration (days) & 360.175 & 282.948 & 229.044 \\
Std. duration (days) & 494.538 & 379.917 & 306.809 \\
\bottomrule
\end{tabular}
\end{table}



\textcolor{uv}{The 9-year gap column refers to a subsample of the Baseline 1992- results; these are the ones that were also active in 2011, but were not in the sample as their completed spells happened before 2011. That is, we select:
\begin{itemize}
    \item people with \textbf{only 1} completed spell in 2011-
    \item Of those, we find out how many would have 2 completed spells if we extended the time window to 1992
    \item We select the ones with a gap between spells of at least 9 years (so effectively those having a very long stable spell in between).
\end{itemize}
The 9 years is because when doing it with 2020, there were 9 years total in the window, so by construction, we are picking people who, if we had the sample time period as in the baseline (10 years) we would not be able to catch them. That is, we would need a time window longer than 9 years to include them. In type lingo, these are likely HL types (high lambda, low delta). The headline result is that the variance is a big larger, but the duration dependence is exactly the same as in the 1992 sample. Circling back to sample selection, this is consistent with the HL types being more prominent in the 1992 sample, as sample selection predicts. I have also added the table just with these three columns.}

\textcolor{uv}{\textbf{Interesting note:} Looking at the sample selection proportions, it checks out: it says we have 49\% of HL types, but we were already capturing 22\% of them in the 10-year sample (2011-), so the new observations we are adding must be 49-22=27\% of the total. The proportion of long gaps over 1992 is 19\%, so roughly consistent.}\\

\begin{table}[htbb!]\label{tab:results_spain_BigGap}
    \centering
    \caption{Spain, large gaps between spells}
\begin{tabular}{llll}
\toprule
Variable & Baseline (2011-) & 9-year gap & Baseline (1992-) \\
\midrule
$\sigma^{2}_y$ & 1.852 & 2.141 & 2.079 \\
$\sigma^{2}_c$ & 1.645 & 1.645 & 1.645 \\
$\sigma^{2}_n$ & 0.098 & 0.336 & 0.001 \\
$\sigma^{2}_e$ & -0.054 & 0.104 & 0.101 \\
$\sigma^{2}_b$ & 0.163 & 0.057 & 0.333 \\
\midrule
\% constant & 0.830 & 0.768 & 0.791 \\
\% spell number & 0.053 & 0.157 & 0.000 \\
\% duration dependence & -0.029 & 0.048 & 0.048 \\
\% heterogeneity & 0.088 & 0.026 & 0.160 \\
\hline \hline
N & 69,406 & 41,604 & 215,912 \\
Average duration (days) & 911.609 & 1750.138 & 927.487 \\
Std. duration (days) & 1045.976 & 1738.254 & 1186.913 \\
\bottomrule
\end{tabular}
\end{table}


\section{Spain: Quits}

\textcolor{uv}{Quits refers to people whose 2 completed spells are quits. Should be combined with the Baseline 2011 column. They have slightly greater variance as the duration dependence switches from negative to positive, but it remains very small. Heterogeneity increases a bit too. These are a very select group, and the number of observations means the estimate is quite imprecise.}

\begin{table}[htbb!]\label{tab:results_spain_Quits}
    \centering
    \caption{Spain, quits}
\begin{tabular}{lll}
\toprule
Variable & Baseline (2011-) & Quits \\
\midrule
$\sigma^{2}_y$ & 1.852 & 2.017 \\
$\sigma^{2}_c$ & 1.645 & 1.645 \\
$\sigma^{2}_n$ & 0.098 & 0.086 \\
$\sigma^{2}_e$ & -0.054 & 0.060 \\
$\sigma^{2}_b$ & 0.163 & 0.226 \\
\midrule
\% constant & 0.830 & 0.815 \\
\% spell number & 0.053 & 0.042 \\
\% duration dependence & -0.029 & 0.030 \\
\% heterogeneity & 0.088 & 0.112 \\
\hline \hline
N & 69,406 & 6,258 \\
Average duration (days) & 911.609 & 831.809 \\
Std. duration (days) & 1045.976 & 1006.983 \\
\bottomrule
\end{tabular}
\end{table}

\section{Descriptive Stats by Spell-bin}

\textcolor{uv}{NOTES: $N_S$ and $N_U$ are number of stable and unstable spells. $T_S$ and $T_U$ are the cumulative duration in years in each kind of spell. For Spain, it checks out: the 1 and 2 bins have an over-representation of the HH types, driving the duration of stable spells down. The 0 bin, in contrast, has very long stable spells and short unstable, reflecting the abundance of HL types. No large differences in observables, except the 0 bin being younger in Spain, and non-nationals dominating in the 0 bin in Germany.}

\textcolor{uv}{As for the short cummualtive duration of stable jobs in the bin 0 in Germany, I asked Claude to exmplain (after checking the code, which is solid) and it said:\\
That's actually explainable from the data itself — let me show you. Looking at the raw CSV for bin 0:\\
  n_PC  mean = 0.18,  SD = 0.38,  \\p10=p25=p50=p75=0,  p90=1\\
  days_p mean = 249,  SD = 759, \\  p10=p25=p50=p75=0,  p90=770\\
  The median of days_p for bin 0 is 0 — more than 75\% of workers in that bin have zero days in stable employment.\\
  The mean is only 0.68 years because 82\% of them have n_PC = 0 (no stable spells at all), and those zeros drag the average down heavily. The handful who do have stable spells (the top ~10\%) are pulling the mean up to 249 days.\\
  This makes sense conceptually: bin 0 workers are people who never completed a full TU cycle. Many of them went into their first TU spell directly (no preceding stable contract), so days_p = 0. The high days_tu (10.65 years) reflects that many are in a very long ongoing/incomplete TU spell — since the spell never closed, it never
  registered as a completed cycle, but all those days still accumulate in days_tu.\\
  So no bug — it's a real feature of the data for that group.}


\begin{table}[]
\caption{Descriptive Stats, conditioning by number of spells}
    \centering
 \begin{tabular}{llllllllll}
\toprule
 &  &  & Age & Female & Non-national & $N_S$ & $N_U$ & $T_S$ & $T_U$ \\
\midrule
Germany & 0 & Mean & 40.63 & 0.48 & 0.23 & 0.51 & 0.00 & 12.79 & 7.37 \\
 &  & Median & 42.00 & 0.00 & 0.00 & 1.00 & 0.00 & 8.35 & 3.72 \\
 &  &  & (9.56) & (0.5) & (0.42) & (0.5) & (0.0) & (12.96) & (8.39) \\
 & 1 & Mean & 40.29 & 0.48 & 0.19 & 1.30 & 1.00 & 12.55 & 7.90 \\
 &  & Median & 40.00 & 0.00 & 0.00 & 1.00 & 1.00 & 10.24 & 5.99 \\
 &  &  & (9.01) & (0.5) & (0.39) & (0.46) & (0.0) & (10.49) & (7.1) \\
 & 2 & Mean & 40.98 & 0.50 & 0.18 & 3.25 & 2.92 & 11.67 & 7.78 \\
 &  & Median & 40.00 & 0.00 & 0.00 & 3.00 & 2.00 & 9.90 & 5.92 \\
 &  &  & (7.93) & (0.5) & (0.39) & (1.5) & (1.44) & (8.5) & (6.54) \\
\midrule
Spain & 0 & Mean & 34.62 & 0.49 & 0.13 & 1.48 & 0.00 & 10.44 & 5.43 \\
 &  & Median & 32.00 & 0.00 & 0.00 & 1.00 & 0.00 & 8.37 & 1.99 \\
 &  &  & (8.09) & (0.5) & (0.33) & (1.61) & (0.0) & (9.8) & (7.5) \\
 & 1 & Mean & 40.45 & 0.47 & 0.14 & 2.06 & 1.00 & 6.99 & 7.55 \\
 &  & Median & 41.00 & 0.00 & 0.00 & 2.00 & 1.00 & 5.48 & 5.65 \\
 &  &  & (8.16) & (0.5) & (0.35) & (1.54) & (0.0) & (6.12) & (6.7) \\
 & 2 & Mean & 41.35 & 0.43 & 0.18 & 3.52 & 2.50 & 6.61 & 7.49 \\
 &  & Median & 41.00 & 0.00 & 0.00 & 3.00 & 2.00 & 5.50 & 6.40 \\
 &  &  & (7.37) & (0.49) & (0.38) & (1.86) & (1.37) & (4.89) & (5.16) \\
\bottomrule
\end{tabular}
    \label{tab:descriptive_stats_bins}
\end{table}

\end{appendices}
\end{document}